% Fiche d'exercices POO Python — classes, attributs, méthodes, encapsulation, héritage (BTS CIEL)
% Compilation :
%   fiche élève : pdflatex Fiche1_POO_Python.tex
%   corrigé     : pdflatex -jobname=Fiche1_POO_Python_Corrige "\def\modecorrige{1}% Fiche d'exercices POO Python — classes, attributs, méthodes, encapsulation, héritage (BTS CIEL)
% Compilation :
%   fiche élève : pdflatex Fiche1_POO_Python.tex
%   corrigé     : pdflatex -jobname=Fiche1_POO_Python_Corrige "\def\modecorrige{1}% Fiche d'exercices POO Python — classes, attributs, méthodes, encapsulation, héritage (BTS CIEL)
% Compilation :
%   fiche élève : pdflatex Fiche1_POO_Python.tex
%   corrigé     : pdflatex -jobname=Fiche1_POO_Python_Corrige "\def\modecorrige{1}% Fiche d'exercices POO Python — classes, attributs, méthodes, encapsulation, héritage (BTS CIEL)
% Compilation :
%   fiche élève : pdflatex Fiche1_POO_Python.tex
%   corrigé     : pdflatex -jobname=Fiche1_POO_Python_Corrige "\def\modecorrige{1}\input{Fiche1_POO_Python.tex}"
\documentclass[10pt,a4paper]{article}

\usepackage[T1]{fontenc}
\usepackage[utf8]{inputenc}
\usepackage{lmodern}
\usepackage[french]{babel}
\usepackage[a4paper,top=0.9cm,bottom=1.1cm,left=1.1cm,right=1.1cm]{geometry}
\usepackage[table]{xcolor}
\usepackage{booktabs,tabularx,array}
\usepackage{enumitem}
\usepackage{needspace}
\usepackage{fancyhdr}
\usepackage{amssymb,amsmath}
\usepackage[most]{tcolorbox}
\usepackage{listings}
\usepackage[strings]{underscore}  % permet d'écrire __init__ sans échappement
\usepackage{hyperref}
\hypersetup{hidelinks}

\providecommand{\modecorrige}{0}
\newif\ifcorrige
\ifnum\modecorrige=1 \corrigetrue \else \corrigefalse \fi

% Palette alignée sur les autres supports du dépôt
\definecolor{ipred}{HTML}{C0392B}
\definecolor{cdcbar}{HTML}{A67C3D}
\definecolor{cdcbg}{HTML}{F4F1EA}
\definecolor{qbar}{HTML}{4A6B5D}
\definecolor{qbg}{HTML}{EDF1EF}
\definecolor{corr}{HTML}{1F5FA8}
\definecolor{corrbg}{HTML}{E8F0FA}
\definecolor{grisfoot}{HTML}{828282}
\definecolor{codebg}{HTML}{F7F7F4}
\definecolor{codekw}{HTML}{8E44AD}
\definecolor{codestr}{HTML}{1F7A4D}
\definecolor{codecom}{HTML}{8A8A8A}

% ---------------------------------------------------------------- code Python
\lstdefinestyle{py}{
  language=Python,
  basicstyle=\ttfamily\scriptsize,
  keywordstyle=\color{codekw}\bfseries,
  stringstyle=\color{codestr},
  commentstyle=\color{codecom}\itshape,
  showstringspaces=false,
  backgroundcolor=\color{codebg},
  frame=single,
  rulecolor=\color{black!25},
  framesep=3pt,
  aboveskip=3pt, belowskip=3pt,
  xleftmargin=2pt, xrightmargin=2pt,
  columns=fullflexible,
  keepspaces=true,
  breaklines=true,
  extendedchars=true,
  literate={é}{{\'e}}1 {è}{{\`e}}1 {ê}{{\^e}}1 {à}{{\`a}}1 {ç}{{\c c}}1
           {ù}{{\`u}}1 {û}{{\^u}}1 {î}{{\^\i}}1 {ô}{{\^o}}1
           {É}{{\'E}}1 {À}{{\`A}}1 {È}{{\`E}}1 {Ç}{{\c C}}1
           {«}{{\guillemotleft}}1 {»}{{\guillemotright}}1,
  morekeywords={self,True,False,None},
}
\lstset{style=py}
\newcommand{\py}[1]{\texttt{#1}}

% --------------------------------------------------------------- encadrés
\newtcolorbox{memo}[1][]{enhanced,breakable,colback=cdcbg,colframe=cdcbar,boxrule=0pt,
  leftrule=2.5pt,arc=0pt,outer arc=0pt,left=5pt,right=5pt,top=3pt,bottom=3pt,
  title={#1},coltitle=cdcbar,fonttitle=\bfseries\footnotesize,attach title to upper={\ :\ },
  fontupper=\footnotesize}

% --------------------------------------------------------------- exercices
\newcounter{exo}
\newcommand{\exo}[1]{%
  \stepcounter{exo}\par\vspace{4pt}
  \ifcorrige\else\noindent\minipage{\textwidth}\fi
  \noindent\textcolor{qbar}{\rule{2.5pt}{10pt}}\hspace{4pt}%
  \textbf{\color{qbar}Exercice \theexo}\ --- \textbf{#1}%
  \par\nobreak\vspace{-3pt}\noindent\textcolor{qbar!35}{\rule{\textwidth}{0.6pt}}%
  \par\nobreak\vspace{1pt}}

% Mise en page de l'énoncé :
%   fiche élève  : code à gauche, questions à droite (gain de place vertical)
%   corrigé      : pleine largeur, l'un sous l'autre (les réponses prennent la place)
\newenvironment{codebloc}
  {\ifcorrige\else\noindent\minipage[t]{0.50\textwidth}\vspace{0pt}\fi}
  {\ifcorrige\else\endminipage\hfill\fi}
\newenvironment{qbloc}
  {\ifcorrige\else\minipage[t]{0.475\textwidth}\vspace{0pt}\fi
   \textbf{\small Questions}\begin{enumerate}[leftmargin=*,nosep,label=\arabic*.]}
  {\end{enumerate}\ifcorrige\else\endminipage\par\fi}

% Variante pleine largeur : exercices sans bloc de code à afficher
\newenvironment{qbloclarge}
  {\par\vspace{2pt}\noindent\textbf{\small Questions}\par\nobreak\begin{enumerate}[leftmargin=*,nosep,label=\arabic*.]}
  {\end{enumerate}}

% Réponse d'une question : rien sur la fiche, texte bleu dans le corrigé
\newcommand{\rep}[1]{\ifcorrige\par\smallskip{\color{corr}\small#1}\par\smallskip\fi}

% Cadre de réponse : présent sur la fiche uniquement.
% \zoneadj retranche la même hauteur à tous les cadres : un seul réglage
% pour ajuster la fiche au nombre de pages voulu.
\newlength{\zoneadj}\setlength{\zoneadj}{0pt}
\newcommand{\zone}[1][4cm]{%
  \ifcorrige\else
    \par\vspace{2pt}\noindent
    \begin{tcolorbox}[colback=white,colframe=black!45,boxrule=0.5pt,arc=1pt,height={\dimexpr#1-\zoneadj\relax},
      left=4pt,right=4pt,top=2pt,bottom=2pt,
      title={\scriptsize\itshape Réponses},coltitle=black!55,fonttitle=\itshape,
      detach title,before upper={\tcbtitle\par}]\end{tcolorbox}%
    \endminipage\par
  \fi}

% Bloc visible seulement dans le corrigé
\newenvironment{solution}
  {\ifcorrige
     \begin{tcolorbox}[enhanced,breakable,colback=corrbg,colframe=corr,boxrule=0.6pt,arc=1pt,
       left=5pt,right=5pt,top=3pt,bottom=3pt,fontupper=\small,
       title={Corrigé},coltitle=corr,fonttitle=\bfseries\small,detach title,
       before upper={\tcbtitle\par\smallskip}]%
   \else
     \setbox0\vbox\bgroup
   \fi}
  {\ifcorrige\end{tcolorbox}\else\egroup\fi}

% Partie
\newcommand{\partie}[1]{%
  \Needspace{7cm}\par\medskip\noindent{\large\bfseries\color{qbar}#1}%
  \par\nobreak\vspace{-2pt}\noindent\textcolor{qbar}{\rule{\textwidth}{1pt}}\par\nobreak\smallskip}

\pagestyle{fancy}
\fancyhf{}
\renewcommand{\headrulewidth}{0pt}
\fancyfoot[L]{\scriptsize\itshape\color{grisfoot}Fiche d'exercices --- POO Python\ifcorrige\ --- CORRIGÉ\fi}
\fancyfoot[R]{\scriptsize\color{grisfoot}\thepage}

\setlist{itemsep=1pt,topsep=2pt,parsep=0pt}
\setlength{\parindent}{0pt}
\setlength{\parskip}{2pt}
\setlength{\columnsep}{8mm}

\begin{document}

\noindent
\begin{minipage}[c]{0.70\textwidth}
{\large\bfseries Fiche d'exercices --- Programmation orientée objet en Python}\\[1pt]
{\small Classes, attributs, méthodes \ $\cdot$\ Encapsulation et héritage \ $\cdot$\ BTS CIEL}
\ifcorrige\\[1pt]{\small\color{corr}\bfseries Version enseignant --- corrigés intégrés}\fi
\end{minipage}\hfill
\begin{minipage}[c]{0.28\textwidth}\small
\ifcorrige\color{corr}corrigé\else Nom :\dotfill\\[3pt] Classe :\dotfill\fi
\end{minipage}

\vspace{2pt}
\begin{memo}[Mémo]
\textbf{Classe} : modèle. \textbf{Objet} (instance) : réalisation concrète de ce modèle.
\textbf{Attribut} : donnée portée par l'objet. \textbf{Méthode} : fonction définie dans la classe.
\py{self} désigne l'objet sur lequel la méthode est appelée ; \py{__init__} est le
\textbf{constructeur}. \py{return} \emph{retourne} une valeur, \py{print} l'\emph{affiche}.
\end{memo}

\partie{Partie A --- Classes, attributs et méthodes}

% ---------------------------------------------------------------- A1
\exo{La classe \py{Voiture}}
\begin{codebloc}
\begin{lstlisting}
class Voiture:
    def __init__(self, marque, couleur):
        self.marque = marque
        self.couleur = couleur

    def afficher_details(self):
        return f"Voiture {self.marque} de couleur {self.couleur}"
\end{lstlisting}
\end{codebloc}
\begin{qbloc}
  \item Que fait la méthode \py{__init__} ?
    \rep{C'est le \textbf{constructeur}. Il est appelé automatiquement à la création de l'objet et
    initialise les attributs d'instance \py{self.marque} et \py{self.couleur} à partir des valeurs
    passées en paramètres.}
  \item Combien d'attributs possède la classe \py{Voiture} ?
    \rep{Deux attributs d'instance : \py{marque} et \py{couleur}.}
  \item Que retourne la méthode \py{afficher_details} ?
    \rep{Une chaîne de caractères (\py{str}) construite par une f-string, par exemple
    \py{"Voiture Toyota de couleur Bleu"}. Elle \emph{retourne} la chaîne, elle ne l'affiche pas.}
  \item Créez une instance de \py{Voiture} de marque \py{"Toyota"} et de couleur \py{"Bleu"},
    appelez \py{afficher_details} et affichez le résultat.
\end{qbloc}
\zone[3.8cm]
\begin{solution}
\begin{lstlisting}
ma_voiture = Voiture("Toyota", "Bleu")
print(ma_voiture.afficher_details())      # Voiture Toyota de couleur Bleu
\end{lstlisting}
\py{print} est nécessaire : sans lui, la chaîne est retournée puis perdue.
\end{solution}

% ---------------------------------------------------------------- A2
\exo{Compléter la classe \py{Chien}}
\begin{codebloc}
\begin{lstlisting}
class Chien:
    def __init__(self, nom, age):
        self.nom = nom
        self.age = age

    # Méthode aboyer qui retourne "Wouf !"
\end{lstlisting}
\end{codebloc}
\begin{qbloc}
  \item Quels sont les attributs de la classe \py{Chien} ?
    \rep{\py{nom} et \py{age}, deux attributs d'instance.}
  \item Comment appelle-t-on la méthode \py{__init__} dans le contexte de la POO ?
    \rep{Le \textbf{constructeur} (ou initialiseur). C'est une \emph{méthode spéciale} ; on ne
    l'appelle jamais directement, elle est déclenchée par \py{Chien("Rex", 3)}.}
  \item Ajoutez une méthode \py{aboyer} qui retourne \py{"Wouf !"}.
  \item Créez un objet \py{Chien} nommé \py{"Rex"}, âgé de 3 ans, et faites-le aboyer.
\end{qbloc}
\zone[3.8cm]
\begin{solution}
\begin{lstlisting}
class Chien:
    def __init__(self, nom, age):
        self.nom = nom
        self.age = age

    def aboyer(self):
        return "Wouf !"

rex = Chien("Rex", 3)
print(rex.aboyer())        # Wouf !
\end{lstlisting}
Le paramètre \py{self} est obligatoire dans la définition, même si \py{aboyer} n'utilise aucun
attribut : il reçoit l'objet sur lequel la méthode est appelée.
\end{solution}

% ---------------------------------------------------------------- A3
\exo{La classe \py{Livre}}
\begin{codebloc}
\begin{lstlisting}
class Livre:
    def __init__(self, titre, auteur):
        self.titre = titre
        self.auteur = auteur

    def description(self):
        return f"{self.titre} par {self.auteur}"
\end{lstlisting}
\end{codebloc}
\begin{qbloc}
  \item Que fait la méthode \py{description} ?
    \rep{Elle retourne une chaîne qui associe le titre et l'auteur du livre, sous la forme
    \py{"titre par auteur"}.}
  \item Combien de paramètres la méthode \py{__init__} attend-elle, en incluant \py{self} ?
    \rep{Trois : \py{self}, \py{titre} et \py{auteur}. À l'appel, seuls deux arguments sont
    fournis : \py{self} est transmis automatiquement par Python.}
  \item Créez un objet \py{Livre} de titre \py{"Le Petit Prince"} et d'auteur
    \py{"Antoine de Saint-Exupéry"}, puis affichez sa description.
\end{qbloc}
\zone[3.8cm]
\begin{solution}
\begin{lstlisting}
livre = Livre("Le Petit Prince", "Antoine de Saint-Exupery")
print(livre.description())      # Le Petit Prince par Antoine de Saint-Exupery
\end{lstlisting}
\end{solution}

% ---------------------------------------------------------------- A4
\exo{La classe \py{Rectangle}}
\begin{codebloc}
\begin{lstlisting}
class Rectangle:
    def __init__(self, longueur, largeur):
        self.longueur = longueur
        self.largeur = largeur

    # Méthode calculer_aire qui retourne l'aire
\end{lstlisting}
\end{codebloc}
\begin{qbloc}
  \item Quels sont les attributs de la classe \py{Rectangle} ?
    \rep{\py{longueur} et \py{largeur}.}
  \item Quelle formule mathématique devez-vous utiliser pour calculer l'aire ?
    \rep{$\text{aire} = \text{longueur} \times \text{largeur}$.}
  \item Ajoutez une méthode \py{calculer_aire} qui retourne l'aire du rectangle.
  \item Créez un rectangle de longueur 5 et de largeur 3, puis affichez son aire.
\end{qbloc}
\zone[8.2cm]
\begin{solution}
\begin{lstlisting}
class Rectangle:
    def __init__(self, longueur, largeur):
        self.longueur = longueur
        self.largeur = largeur

    def calculer_aire(self):
        return self.longueur * self.largeur

r = Rectangle(5, 3)
print(r.calculer_aire())        # 15
\end{lstlisting}
Une méthode peut aussi être ajoutée pour le périmètre :\\
\py{return 2 * (self.longueur + self.largeur)}.
\end{solution}

% ---------------------------------------------------------------- A5
\exo{Création complète d'une classe \py{Etudiant}}
On souhaite modéliser un étudiant par son nom et sa note.
\begin{qbloclarge}
  \item Créez une classe \py{Etudiant} avec un constructeur \py{__init__} prenant en paramètres
    \py{nom} et \py{note}, et deux attributs d'instance \py{self.nom} et \py{self.note}.
  \item Ajoutez une méthode \py{a_reussi} qui retourne \py{True} si la note est supérieure ou égale
    à 10, et \py{False} sinon.
  \item Testez votre classe en créant un étudiant avec une note de 8.
  \item Ajoutez une méthode \py{augmenter_note} qui augmente la note de 2 points.
  \item Créez un étudiant nommé \py{"Alice"} avec une note de 15 et affichez si elle a réussi,
    grâce à la méthode \py{a_reussi}.
\end{qbloclarge}
\zone[8.7cm]
\begin{solution}
\begin{lstlisting}
class Etudiant:
    def __init__(self, nom, note):
        self.nom = nom
        self.note = note

    def a_reussi(self):
        return self.note >= 10

    def augmenter_note(self):
        self.note = self.note + 2

e = Etudiant("Bob", 8)          # 3. test avec une note de 8
print(e.a_reussi())             # False

alice = Etudiant("Alice", 15)   # 5. Alice
print(alice.a_reussi())         # True
\end{lstlisting}
\textbf{Points d'attention.} \py{return self.note >= 10} suffit : la comparaison produit déjà un
booléen, un \py{if ... return True else return False} est inutile. \py{augmenter_note} \emph{modifie}
l'attribut et ne retourne rien : elle agit sur l'état de l'objet.
\end{solution}

\partie{Partie B --- Encapsulation et héritage}

\begin{memo}[Mémo]
\textbf{Encapsulation} : \py{self.__nom} (deux soulignés) déclenche le \emph{name mangling} — Python
renomme l'attribut en \py{_Classe__nom}, donc inaccessible de l'extérieur par son nom d'origine. On y
accède par un \textbf{accesseur} (\py{get_nom}) et on le modifie par un \textbf{mutateur}
(\py{set_nom}), qui peut valider la valeur.
\textbf{Héritage} : \py{class Chat(Animal)} — la classe fille reçoit les attributs et méthodes de la
mère, peut en ajouter ou \textbf{redéfinir} les siennes ; \py{super().__init__(...)} appelle le
constructeur de la mère.
\end{memo}

% ---------------------------------------------------------------- B1
\exo{Encapsulation simple avec attribut privé}
\begin{codebloc}
\begin{lstlisting}
class Personne:
    def __init__(self, nom, age):
        self.__nom = nom
        self.__age = age

    def get_nom(self):
        return self.__nom

    # Méthode get_age()
\end{lstlisting}
\end{codebloc}
\begin{qbloc}
  \item Pourquoi les attributs \py{__nom} et \py{__age} sont-ils privés ?
    \rep{Pour empêcher leur modification directe depuis l'extérieur. L'accès passe par des méthodes,
    ce qui permet de contrôler la lecture et l'écriture et de modifier l'implémentation interne sans
    casser le code qui utilise la classe.}
  \item Comment accède-t-on à l'attribut \py{__nom} ?
    \rep{Par l'accesseur : \py{p.get_nom()}. L'écriture \py{p.__nom} provoque une
    \py{AttributeError}.}
  \item Ajoutez une méthode \py{get_age} pour accéder à l'attribut \py{__age}.
  \item Créez une instance de \py{Personne} de nom \py{"Luc"} et d'âge 25, puis affichez son nom.
\end{qbloc}
\zone[4.0cm]
\begin{solution}
\begin{lstlisting}
    def get_age(self):
        return self.__age

p = Personne("Luc", 25)
print(p.get_nom())        # Luc
print(p.get_age())        # 25
\end{lstlisting}
\end{solution}

% ---------------------------------------------------------------- B2
\exo{Ajouter un mutateur}
\begin{codebloc}
\begin{lstlisting}
class CompteBancaire:
    def __init__(self, titulaire, solde):
        self.__titulaire = titulaire
        self.__solde = solde

    def get_solde(self):
        return self.__solde
\end{lstlisting}
\end{codebloc}
\begin{qbloc}
  \item Que signifie l'encapsulation dans cette classe ?
    \rep{Le titulaire et le solde sont privés : aucun code extérieur ne peut écrire
    \py{compte.__solde = 1000000}. Le solde ne peut être consulté que par \py{get_solde} et modifié
    que par les méthodes que la classe met à disposition.}
  \item Que fait la méthode \py{get_solde} ?
    \rep{Elle retourne la valeur de l'attribut privé \py{__solde} : c'est un accesseur en lecture
    seule.}
  \item Créez un compte bancaire pour \py{"Emma"} avec un solde de 500 et affichez le solde.
  \item Ajoutez une méthode \py{set_solde} pour modifier le solde.
\end{qbloc}
\zone[4.0cm]
\begin{solution}
\begin{lstlisting}
    def set_solde(self, nouveau_solde):
        self.__solde = nouveau_solde

compte = CompteBancaire("Emma", 500)
print(compte.get_solde())       # 500
compte.set_solde(750)
print(compte.get_solde())       # 750
\end{lstlisting}
\textbf{Prolongement.} Un mutateur brut n'apporte aucune sécurité. Une version plus réaliste refuse
un solde négatif, ou expose plutôt des opérations métier : \py{deposer(montant)} et
\py{retirer(montant)} avec vérification de la provision.
\end{solution}

% ---------------------------------------------------------------- B3
\exo{Encapsulation avec validation}
\begin{codebloc}
\begin{lstlisting}
class Etudiant:
    def __init__(self, nom, note):
        self.__nom = nom
        self.__note = note

    def get_note(self):
        return self.__note

    # Méthode set_note : notes entre 0 et 20
\end{lstlisting}
\end{codebloc}
\begin{qbloc}
  \item Pourquoi utiliser un mutateur pour la note ?
    \rep{Pour garantir un \emph{invariant} : une note reste comprise entre 0 et 20. Le mutateur est
    le seul point d'entrée en écriture, il peut donc vérifier la valeur avant de l'enregistrer.}
  \item Que se passe-t-il si on essaie d'accéder directement à \py{__note} ?
    \rep{Python lève une \py{AttributeError} : à cause du \emph{name mangling}, l'attribut s'appelle
    en réalité \py{_Etudiant__note} et le nom \py{__note} n'existe pas depuis l'extérieur.}
  \item Ajoutez une méthode \py{set_note} qui n'accepte que des notes entre 0 et 20.
  \item Créez un étudiant nommé \py{"Paul"} avec une note de 15, puis essayez de changer la note
    à 25.
\end{qbloc}
\zone[4.0cm]
\begin{solution}
\begin{lstlisting}
    def set_note(self, nouvelle_note):
        if 0 <= nouvelle_note <= 20:
            self.__note = nouvelle_note
        else:
            print("Note invalide : elle doit etre comprise entre 0 et 20.")

paul = Etudiant("Paul", 15)
paul.set_note(25)         # Note invalide : ...
print(paul.get_note())    # 15 : la note n'a pas change
\end{lstlisting}
\textbf{Variante professionnelle.} Au lieu d'un \py{print}, lever une exception rend l'erreur
impossible à ignorer :\\
\py{raise ValueError("note hors intervalle")}.
\end{solution}

% ---------------------------------------------------------------- B4
\exo{Héritage simple}
\begin{codebloc}
\begin{lstlisting}
class Animal:
    def __init__(self, nom):
        self.__nom = nom

    def get_nom(self):
        return self.__nom

    def faire_son(self):
        return "Son generique"

class Chat(Animal):
    # Méthode faire_son pour retourner "Miaou"
\end{lstlisting}
\end{codebloc}
\begin{qbloc}
  \item Quelle classe est la classe parent ?
    \rep{\py{Animal}. \py{Chat} est la classe fille (ou classe dérivée).}
  \item Que fait l'héritage dans ce cas ?
    \rep{\py{Chat} récupère sans les réécrire le constructeur \py{__init__}, l'attribut privé
    \py{__nom} et la méthode \py{get_nom} d'\py{Animal}. Elle peut en plus redéfinir \py{faire_son}
    pour lui donner son comportement propre.}
  \item Redéfinissez la méthode \py{faire_son} pour retourner \py{"Miaou"}.
  \item Créez un \py{Chat} nommé \py{"Minou"} et affichez son son.
\end{qbloc}
\zone[4.0cm]
\begin{solution}
\begin{lstlisting}
class Chat(Animal):
    def faire_son(self):
        return "Miaou"

minou = Chat("Minou")
print(minou.get_nom())     # Minou  (methode heritee)
print(minou.faire_son())   # Miaou  (methode redefinie)
\end{lstlisting}
\py{Chat} ne définit pas de \py{__init__} : celui d'\py{Animal} est utilisé tel quel. Appeler
\py{faire_son()} sur un \py{Chat} exécute la version de \py{Chat} : c'est le
\textbf{polymorphisme}.
\end{solution}

% ---------------------------------------------------------------- B5
\exo{Héritage avec nouvel attribut}
\begin{codebloc}
\begin{lstlisting}
class Vehicule:
    def __init__(self, marque):
        self.__marque = marque

    def get_marque(self):
        return self.__marque

class Voiture(Vehicule):
    def __init__(self, marque, couleur):
        super().__init__(marque)
        self.__couleur = couleur

    # À compléter : méthode get_couleur
\end{lstlisting}
\end{codebloc}
\begin{qbloc}
  \item Que fait \py{super().__init__(marque)} ?
    \rep{Il appelle le constructeur de la classe mère \py{Vehicule} pour initialiser la partie
    héritée de l'objet, ici l'attribut privé \py{__marque}. Sans cet appel, \py{get_marque}
    échouerait car l'attribut n'aurait jamais été créé.}
  \item Combien d'attributs privés possède la classe \py{Voiture} ?
    \rep{Deux au total dans l'objet : \py{_Vehicule__marque}, hérité et initialisé par \py{super()},
    et \py{_Voiture__couleur}, propre à \py{Voiture}. La classe \py{Voiture} n'en déclare qu'un seul
    elle-même.}
  \item À compléter : ajoutez une méthode \py{get_couleur}.
\end{qbloc}
\zone[4.0cm]
\begin{solution}
\begin{lstlisting}
    def get_couleur(self):
        return self.__couleur

v = Voiture("Renault", "Bleu")
print(v.get_marque())      # Renault
print(v.get_couleur())     # Bleu
\end{lstlisting}
\textcolor{ipred}{\textbf{Piège.}} \py{__marque} est privé \emph{à \py{Vehicule}} : à cause du
\emph{name mangling}, une méthode de \py{Voiture} ne peut pas écrire \py{self.__marque}. Elle doit
passer par \py{self.get_marque()}. C'est exactement le rôle de l'accesseur hérité.
\end{solution}

% ---------------------------------------------------------------- B6
\exo{Encapsulation et mutateur avec validation}
\begin{codebloc}
\begin{lstlisting}
class Livre:
    def __init__(self, titre, prix):
        self.__titre = titre
        self.__prix = prix

    def get_prix(self):
        return self.__prix

    # À compléter : set_prix refuse les prix négatifs
\end{lstlisting}
\end{codebloc}
\begin{qbloc}
  \item Pourquoi le prix est-il un attribut privé ?
    \rep{Parce qu'un prix doit respecter une règle métier : il ne peut pas être négatif. Le rendre
    privé interdit l'écriture directe et oblige à passer par une méthode qui applique cette règle.}
  \item Que fait un mutateur avec validation ?
    \rep{Il vérifie la valeur reçue avant de l'affecter. Si elle est valide, il met l'attribut à
    jour ; sinon il refuse la modification et l'objet conserve son état précédent, qui reste
    cohérent.}
  \item Créez un \py{Livre} intitulé \py{"Python 101"} avec un prix de 30, puis essayez de définir
    un prix de $-10$.
\end{qbloc}
\zone[4.0cm]
\begin{solution}
\begin{lstlisting}
    def set_prix(self, nouveau_prix):
        if nouveau_prix >= 0:
            self.__prix = nouveau_prix
        else:
            print("Prix invalide : un prix ne peut pas etre negatif.")

livre = Livre("Python 101", 30)
livre.set_prix(-10)        # Prix invalide : ...
print(livre.get_prix())    # 30 : le prix est inchange
\end{lstlisting}
\textbf{Remarque.} Le constructeur devrait appliquer la même validation, par exemple en appelant
\py{self.set_prix(prix)} au lieu d'affecter directement : sinon \py{Livre("X", -10)} crée un objet
invalide dès le départ.
\end{solution}

% ============================================================ synthèse (corrigé)
\ifcorrige
\clearpage
\partie{Synthèse}
\begin{center}\small
\begin{tabularx}{\textwidth}{@{}>{\bfseries}p{4cm} X@{}}
\toprule
Notion & Ce qu'il faut retenir \\
\midrule
Classe & Modèle décrivant des données et des comportements : \py{class Voiture:} \\
Objet / instance & Réalisation concrète de la classe : \py{v = Voiture("Toyota", "Bleu")} \\
Constructeur & \py{__init__}, appelé automatiquement, initialise les attributs \\
\py{self} & Référence à l'objet courant ; premier paramètre de toute méthode d'instance \\
Attribut d'instance & Donnée propre à l'objet : \py{self.marque} \\
Méthode & Fonction définie dans la classe, appelée par \py{objet.methode()} \\
Attribut privé & \py{self.__prix} : inaccessible depuis l'extérieur (\emph{name mangling}) \\
Accesseur (\emph{getter}) & \py{get_prix()} : lit un attribut privé \\
Mutateur (\emph{setter}) & \py{set_prix(v)} : écrit un attribut privé après validation \\
Héritage & \py{class Voiture(Vehicule)} : la fille reçoit les membres de la mère \\
\py{super()} & Accède à la classe mère, notamment \py{super().__init__(...)} \\
Redéfinition & La fille réécrit une méthode de la mère (\py{faire_son}) \\
\bottomrule
\end{tabularx}
\end{center}

\begin{tcolorbox}[enhanced,breakable,colback=red!4,colframe=ipred,boxrule=0pt,leftrule=2.5pt,
  arc=0pt,outer arc=0pt,left=5pt,right=5pt,top=3pt,bottom=3pt,fontupper=\small]
\textbf{Erreurs les plus fréquentes.} Oublier \py{self} dans la définition d'une méthode ;
confondre \py{return} (retourne une valeur) et \py{print} (affiche) ; écrire \py{objet.__attribut}
depuis l'extérieur ; oublier \py{super().__init__(...)} dans le constructeur d'une classe fille ;
appeler une méthode sans les parenthèses (\py{v.afficher_details} retourne la méthode, pas son
résultat).
\end{tcolorbox}
\fi

\end{document}
"
\documentclass[10pt,a4paper]{article}

\usepackage[T1]{fontenc}
\usepackage[utf8]{inputenc}
\usepackage{lmodern}
\usepackage[french]{babel}
\usepackage[a4paper,top=0.9cm,bottom=1.1cm,left=1.1cm,right=1.1cm]{geometry}
\usepackage[table]{xcolor}
\usepackage{booktabs,tabularx,array}
\usepackage{enumitem}
\usepackage{needspace}
\usepackage{fancyhdr}
\usepackage{amssymb,amsmath}
\usepackage[most]{tcolorbox}
\usepackage{listings}
\usepackage[strings]{underscore}  % permet d'écrire __init__ sans échappement
\usepackage{hyperref}
\hypersetup{hidelinks}

\providecommand{\modecorrige}{0}
\newif\ifcorrige
\ifnum\modecorrige=1 \corrigetrue \else \corrigefalse \fi

% Palette alignée sur les autres supports du dépôt
\definecolor{ipred}{HTML}{C0392B}
\definecolor{cdcbar}{HTML}{A67C3D}
\definecolor{cdcbg}{HTML}{F4F1EA}
\definecolor{qbar}{HTML}{4A6B5D}
\definecolor{qbg}{HTML}{EDF1EF}
\definecolor{corr}{HTML}{1F5FA8}
\definecolor{corrbg}{HTML}{E8F0FA}
\definecolor{grisfoot}{HTML}{828282}
\definecolor{codebg}{HTML}{F7F7F4}
\definecolor{codekw}{HTML}{8E44AD}
\definecolor{codestr}{HTML}{1F7A4D}
\definecolor{codecom}{HTML}{8A8A8A}

% ---------------------------------------------------------------- code Python
\lstdefinestyle{py}{
  language=Python,
  basicstyle=\ttfamily\scriptsize,
  keywordstyle=\color{codekw}\bfseries,
  stringstyle=\color{codestr},
  commentstyle=\color{codecom}\itshape,
  showstringspaces=false,
  backgroundcolor=\color{codebg},
  frame=single,
  rulecolor=\color{black!25},
  framesep=3pt,
  aboveskip=3pt, belowskip=3pt,
  xleftmargin=2pt, xrightmargin=2pt,
  columns=fullflexible,
  keepspaces=true,
  breaklines=true,
  extendedchars=true,
  literate={é}{{\'e}}1 {è}{{\`e}}1 {ê}{{\^e}}1 {à}{{\`a}}1 {ç}{{\c c}}1
           {ù}{{\`u}}1 {û}{{\^u}}1 {î}{{\^\i}}1 {ô}{{\^o}}1
           {É}{{\'E}}1 {À}{{\`A}}1 {È}{{\`E}}1 {Ç}{{\c C}}1
           {«}{{\guillemotleft}}1 {»}{{\guillemotright}}1,
  morekeywords={self,True,False,None},
}
\lstset{style=py}
\newcommand{\py}[1]{\texttt{#1}}

% --------------------------------------------------------------- encadrés
\newtcolorbox{memo}[1][]{enhanced,breakable,colback=cdcbg,colframe=cdcbar,boxrule=0pt,
  leftrule=2.5pt,arc=0pt,outer arc=0pt,left=5pt,right=5pt,top=3pt,bottom=3pt,
  title={#1},coltitle=cdcbar,fonttitle=\bfseries\footnotesize,attach title to upper={\ :\ },
  fontupper=\footnotesize}

% --------------------------------------------------------------- exercices
\newcounter{exo}
\newcommand{\exo}[1]{%
  \stepcounter{exo}\par\vspace{4pt}
  \ifcorrige\else\noindent\minipage{\textwidth}\fi
  \noindent\textcolor{qbar}{\rule{2.5pt}{10pt}}\hspace{4pt}%
  \textbf{\color{qbar}Exercice \theexo}\ --- \textbf{#1}%
  \par\nobreak\vspace{-3pt}\noindent\textcolor{qbar!35}{\rule{\textwidth}{0.6pt}}%
  \par\nobreak\vspace{1pt}}

% Mise en page de l'énoncé :
%   fiche élève  : code à gauche, questions à droite (gain de place vertical)
%   corrigé      : pleine largeur, l'un sous l'autre (les réponses prennent la place)
\newenvironment{codebloc}
  {\ifcorrige\else\noindent\minipage[t]{0.50\textwidth}\vspace{0pt}\fi}
  {\ifcorrige\else\endminipage\hfill\fi}
\newenvironment{qbloc}
  {\ifcorrige\else\minipage[t]{0.475\textwidth}\vspace{0pt}\fi
   \textbf{\small Questions}\begin{enumerate}[leftmargin=*,nosep,label=\arabic*.]}
  {\end{enumerate}\ifcorrige\else\endminipage\par\fi}

% Variante pleine largeur : exercices sans bloc de code à afficher
\newenvironment{qbloclarge}
  {\par\vspace{2pt}\noindent\textbf{\small Questions}\par\nobreak\begin{enumerate}[leftmargin=*,nosep,label=\arabic*.]}
  {\end{enumerate}}

% Réponse d'une question : rien sur la fiche, texte bleu dans le corrigé
\newcommand{\rep}[1]{\ifcorrige\par\smallskip{\color{corr}\small#1}\par\smallskip\fi}

% Cadre de réponse : présent sur la fiche uniquement.
% \zoneadj retranche la même hauteur à tous les cadres : un seul réglage
% pour ajuster la fiche au nombre de pages voulu.
\newlength{\zoneadj}\setlength{\zoneadj}{0pt}
\newcommand{\zone}[1][4cm]{%
  \ifcorrige\else
    \par\vspace{2pt}\noindent
    \begin{tcolorbox}[colback=white,colframe=black!45,boxrule=0.5pt,arc=1pt,height={\dimexpr#1-\zoneadj\relax},
      left=4pt,right=4pt,top=2pt,bottom=2pt,
      title={\scriptsize\itshape Réponses},coltitle=black!55,fonttitle=\itshape,
      detach title,before upper={\tcbtitle\par}]\end{tcolorbox}%
    \endminipage\par
  \fi}

% Bloc visible seulement dans le corrigé
\newenvironment{solution}
  {\ifcorrige
     \begin{tcolorbox}[enhanced,breakable,colback=corrbg,colframe=corr,boxrule=0.6pt,arc=1pt,
       left=5pt,right=5pt,top=3pt,bottom=3pt,fontupper=\small,
       title={Corrigé},coltitle=corr,fonttitle=\bfseries\small,detach title,
       before upper={\tcbtitle\par\smallskip}]%
   \else
     \setbox0\vbox\bgroup
   \fi}
  {\ifcorrige\end{tcolorbox}\else\egroup\fi}

% Partie
\newcommand{\partie}[1]{%
  \Needspace{7cm}\par\medskip\noindent{\large\bfseries\color{qbar}#1}%
  \par\nobreak\vspace{-2pt}\noindent\textcolor{qbar}{\rule{\textwidth}{1pt}}\par\nobreak\smallskip}

\pagestyle{fancy}
\fancyhf{}
\renewcommand{\headrulewidth}{0pt}
\fancyfoot[L]{\scriptsize\itshape\color{grisfoot}Fiche d'exercices --- POO Python\ifcorrige\ --- CORRIGÉ\fi}
\fancyfoot[R]{\scriptsize\color{grisfoot}\thepage}

\setlist{itemsep=1pt,topsep=2pt,parsep=0pt}
\setlength{\parindent}{0pt}
\setlength{\parskip}{2pt}
\setlength{\columnsep}{8mm}

\begin{document}

\noindent
\begin{minipage}[c]{0.70\textwidth}
{\large\bfseries Fiche d'exercices --- Programmation orientée objet en Python}\\[1pt]
{\small Classes, attributs, méthodes \ $\cdot$\ Encapsulation et héritage \ $\cdot$\ BTS CIEL}
\ifcorrige\\[1pt]{\small\color{corr}\bfseries Version enseignant --- corrigés intégrés}\fi
\end{minipage}\hfill
\begin{minipage}[c]{0.28\textwidth}\small
\ifcorrige\color{corr}corrigé\else Nom :\dotfill\\[3pt] Classe :\dotfill\fi
\end{minipage}

\vspace{2pt}
\begin{memo}[Mémo]
\textbf{Classe} : modèle. \textbf{Objet} (instance) : réalisation concrète de ce modèle.
\textbf{Attribut} : donnée portée par l'objet. \textbf{Méthode} : fonction définie dans la classe.
\py{self} désigne l'objet sur lequel la méthode est appelée ; \py{__init__} est le
\textbf{constructeur}. \py{return} \emph{retourne} une valeur, \py{print} l'\emph{affiche}.
\end{memo}

\partie{Partie A --- Classes, attributs et méthodes}

% ---------------------------------------------------------------- A1
\exo{La classe \py{Voiture}}
\begin{codebloc}
\begin{lstlisting}
class Voiture:
    def __init__(self, marque, couleur):
        self.marque = marque
        self.couleur = couleur

    def afficher_details(self):
        return f"Voiture {self.marque} de couleur {self.couleur}"
\end{lstlisting}
\end{codebloc}
\begin{qbloc}
  \item Que fait la méthode \py{__init__} ?
    \rep{C'est le \textbf{constructeur}. Il est appelé automatiquement à la création de l'objet et
    initialise les attributs d'instance \py{self.marque} et \py{self.couleur} à partir des valeurs
    passées en paramètres.}
  \item Combien d'attributs possède la classe \py{Voiture} ?
    \rep{Deux attributs d'instance : \py{marque} et \py{couleur}.}
  \item Que retourne la méthode \py{afficher_details} ?
    \rep{Une chaîne de caractères (\py{str}) construite par une f-string, par exemple
    \py{"Voiture Toyota de couleur Bleu"}. Elle \emph{retourne} la chaîne, elle ne l'affiche pas.}
  \item Créez une instance de \py{Voiture} de marque \py{"Toyota"} et de couleur \py{"Bleu"},
    appelez \py{afficher_details} et affichez le résultat.
\end{qbloc}
\zone[3.8cm]
\begin{solution}
\begin{lstlisting}
ma_voiture = Voiture("Toyota", "Bleu")
print(ma_voiture.afficher_details())      # Voiture Toyota de couleur Bleu
\end{lstlisting}
\py{print} est nécessaire : sans lui, la chaîne est retournée puis perdue.
\end{solution}

% ---------------------------------------------------------------- A2
\exo{Compléter la classe \py{Chien}}
\begin{codebloc}
\begin{lstlisting}
class Chien:
    def __init__(self, nom, age):
        self.nom = nom
        self.age = age

    # Méthode aboyer qui retourne "Wouf !"
\end{lstlisting}
\end{codebloc}
\begin{qbloc}
  \item Quels sont les attributs de la classe \py{Chien} ?
    \rep{\py{nom} et \py{age}, deux attributs d'instance.}
  \item Comment appelle-t-on la méthode \py{__init__} dans le contexte de la POO ?
    \rep{Le \textbf{constructeur} (ou initialiseur). C'est une \emph{méthode spéciale} ; on ne
    l'appelle jamais directement, elle est déclenchée par \py{Chien("Rex", 3)}.}
  \item Ajoutez une méthode \py{aboyer} qui retourne \py{"Wouf !"}.
  \item Créez un objet \py{Chien} nommé \py{"Rex"}, âgé de 3 ans, et faites-le aboyer.
\end{qbloc}
\zone[3.8cm]
\begin{solution}
\begin{lstlisting}
class Chien:
    def __init__(self, nom, age):
        self.nom = nom
        self.age = age

    def aboyer(self):
        return "Wouf !"

rex = Chien("Rex", 3)
print(rex.aboyer())        # Wouf !
\end{lstlisting}
Le paramètre \py{self} est obligatoire dans la définition, même si \py{aboyer} n'utilise aucun
attribut : il reçoit l'objet sur lequel la méthode est appelée.
\end{solution}

% ---------------------------------------------------------------- A3
\exo{La classe \py{Livre}}
\begin{codebloc}
\begin{lstlisting}
class Livre:
    def __init__(self, titre, auteur):
        self.titre = titre
        self.auteur = auteur

    def description(self):
        return f"{self.titre} par {self.auteur}"
\end{lstlisting}
\end{codebloc}
\begin{qbloc}
  \item Que fait la méthode \py{description} ?
    \rep{Elle retourne une chaîne qui associe le titre et l'auteur du livre, sous la forme
    \py{"titre par auteur"}.}
  \item Combien de paramètres la méthode \py{__init__} attend-elle, en incluant \py{self} ?
    \rep{Trois : \py{self}, \py{titre} et \py{auteur}. À l'appel, seuls deux arguments sont
    fournis : \py{self} est transmis automatiquement par Python.}
  \item Créez un objet \py{Livre} de titre \py{"Le Petit Prince"} et d'auteur
    \py{"Antoine de Saint-Exupéry"}, puis affichez sa description.
\end{qbloc}
\zone[3.8cm]
\begin{solution}
\begin{lstlisting}
livre = Livre("Le Petit Prince", "Antoine de Saint-Exupery")
print(livre.description())      # Le Petit Prince par Antoine de Saint-Exupery
\end{lstlisting}
\end{solution}

% ---------------------------------------------------------------- A4
\exo{La classe \py{Rectangle}}
\begin{codebloc}
\begin{lstlisting}
class Rectangle:
    def __init__(self, longueur, largeur):
        self.longueur = longueur
        self.largeur = largeur

    # Méthode calculer_aire qui retourne l'aire
\end{lstlisting}
\end{codebloc}
\begin{qbloc}
  \item Quels sont les attributs de la classe \py{Rectangle} ?
    \rep{\py{longueur} et \py{largeur}.}
  \item Quelle formule mathématique devez-vous utiliser pour calculer l'aire ?
    \rep{$\text{aire} = \text{longueur} \times \text{largeur}$.}
  \item Ajoutez une méthode \py{calculer_aire} qui retourne l'aire du rectangle.
  \item Créez un rectangle de longueur 5 et de largeur 3, puis affichez son aire.
\end{qbloc}
\zone[8.2cm]
\begin{solution}
\begin{lstlisting}
class Rectangle:
    def __init__(self, longueur, largeur):
        self.longueur = longueur
        self.largeur = largeur

    def calculer_aire(self):
        return self.longueur * self.largeur

r = Rectangle(5, 3)
print(r.calculer_aire())        # 15
\end{lstlisting}
Une méthode peut aussi être ajoutée pour le périmètre :\\
\py{return 2 * (self.longueur + self.largeur)}.
\end{solution}

% ---------------------------------------------------------------- A5
\exo{Création complète d'une classe \py{Etudiant}}
On souhaite modéliser un étudiant par son nom et sa note.
\begin{qbloclarge}
  \item Créez une classe \py{Etudiant} avec un constructeur \py{__init__} prenant en paramètres
    \py{nom} et \py{note}, et deux attributs d'instance \py{self.nom} et \py{self.note}.
  \item Ajoutez une méthode \py{a_reussi} qui retourne \py{True} si la note est supérieure ou égale
    à 10, et \py{False} sinon.
  \item Testez votre classe en créant un étudiant avec une note de 8.
  \item Ajoutez une méthode \py{augmenter_note} qui augmente la note de 2 points.
  \item Créez un étudiant nommé \py{"Alice"} avec une note de 15 et affichez si elle a réussi,
    grâce à la méthode \py{a_reussi}.
\end{qbloclarge}
\zone[8.7cm]
\begin{solution}
\begin{lstlisting}
class Etudiant:
    def __init__(self, nom, note):
        self.nom = nom
        self.note = note

    def a_reussi(self):
        return self.note >= 10

    def augmenter_note(self):
        self.note = self.note + 2

e = Etudiant("Bob", 8)          # 3. test avec une note de 8
print(e.a_reussi())             # False

alice = Etudiant("Alice", 15)   # 5. Alice
print(alice.a_reussi())         # True
\end{lstlisting}
\textbf{Points d'attention.} \py{return self.note >= 10} suffit : la comparaison produit déjà un
booléen, un \py{if ... return True else return False} est inutile. \py{augmenter_note} \emph{modifie}
l'attribut et ne retourne rien : elle agit sur l'état de l'objet.
\end{solution}

\partie{Partie B --- Encapsulation et héritage}

\begin{memo}[Mémo]
\textbf{Encapsulation} : \py{self.__nom} (deux soulignés) déclenche le \emph{name mangling} — Python
renomme l'attribut en \py{_Classe__nom}, donc inaccessible de l'extérieur par son nom d'origine. On y
accède par un \textbf{accesseur} (\py{get_nom}) et on le modifie par un \textbf{mutateur}
(\py{set_nom}), qui peut valider la valeur.
\textbf{Héritage} : \py{class Chat(Animal)} — la classe fille reçoit les attributs et méthodes de la
mère, peut en ajouter ou \textbf{redéfinir} les siennes ; \py{super().__init__(...)} appelle le
constructeur de la mère.
\end{memo}

% ---------------------------------------------------------------- B1
\exo{Encapsulation simple avec attribut privé}
\begin{codebloc}
\begin{lstlisting}
class Personne:
    def __init__(self, nom, age):
        self.__nom = nom
        self.__age = age

    def get_nom(self):
        return self.__nom

    # Méthode get_age()
\end{lstlisting}
\end{codebloc}
\begin{qbloc}
  \item Pourquoi les attributs \py{__nom} et \py{__age} sont-ils privés ?
    \rep{Pour empêcher leur modification directe depuis l'extérieur. L'accès passe par des méthodes,
    ce qui permet de contrôler la lecture et l'écriture et de modifier l'implémentation interne sans
    casser le code qui utilise la classe.}
  \item Comment accède-t-on à l'attribut \py{__nom} ?
    \rep{Par l'accesseur : \py{p.get_nom()}. L'écriture \py{p.__nom} provoque une
    \py{AttributeError}.}
  \item Ajoutez une méthode \py{get_age} pour accéder à l'attribut \py{__age}.
  \item Créez une instance de \py{Personne} de nom \py{"Luc"} et d'âge 25, puis affichez son nom.
\end{qbloc}
\zone[4.0cm]
\begin{solution}
\begin{lstlisting}
    def get_age(self):
        return self.__age

p = Personne("Luc", 25)
print(p.get_nom())        # Luc
print(p.get_age())        # 25
\end{lstlisting}
\end{solution}

% ---------------------------------------------------------------- B2
\exo{Ajouter un mutateur}
\begin{codebloc}
\begin{lstlisting}
class CompteBancaire:
    def __init__(self, titulaire, solde):
        self.__titulaire = titulaire
        self.__solde = solde

    def get_solde(self):
        return self.__solde
\end{lstlisting}
\end{codebloc}
\begin{qbloc}
  \item Que signifie l'encapsulation dans cette classe ?
    \rep{Le titulaire et le solde sont privés : aucun code extérieur ne peut écrire
    \py{compte.__solde = 1000000}. Le solde ne peut être consulté que par \py{get_solde} et modifié
    que par les méthodes que la classe met à disposition.}
  \item Que fait la méthode \py{get_solde} ?
    \rep{Elle retourne la valeur de l'attribut privé \py{__solde} : c'est un accesseur en lecture
    seule.}
  \item Créez un compte bancaire pour \py{"Emma"} avec un solde de 500 et affichez le solde.
  \item Ajoutez une méthode \py{set_solde} pour modifier le solde.
\end{qbloc}
\zone[4.0cm]
\begin{solution}
\begin{lstlisting}
    def set_solde(self, nouveau_solde):
        self.__solde = nouveau_solde

compte = CompteBancaire("Emma", 500)
print(compte.get_solde())       # 500
compte.set_solde(750)
print(compte.get_solde())       # 750
\end{lstlisting}
\textbf{Prolongement.} Un mutateur brut n'apporte aucune sécurité. Une version plus réaliste refuse
un solde négatif, ou expose plutôt des opérations métier : \py{deposer(montant)} et
\py{retirer(montant)} avec vérification de la provision.
\end{solution}

% ---------------------------------------------------------------- B3
\exo{Encapsulation avec validation}
\begin{codebloc}
\begin{lstlisting}
class Etudiant:
    def __init__(self, nom, note):
        self.__nom = nom
        self.__note = note

    def get_note(self):
        return self.__note

    # Méthode set_note : notes entre 0 et 20
\end{lstlisting}
\end{codebloc}
\begin{qbloc}
  \item Pourquoi utiliser un mutateur pour la note ?
    \rep{Pour garantir un \emph{invariant} : une note reste comprise entre 0 et 20. Le mutateur est
    le seul point d'entrée en écriture, il peut donc vérifier la valeur avant de l'enregistrer.}
  \item Que se passe-t-il si on essaie d'accéder directement à \py{__note} ?
    \rep{Python lève une \py{AttributeError} : à cause du \emph{name mangling}, l'attribut s'appelle
    en réalité \py{_Etudiant__note} et le nom \py{__note} n'existe pas depuis l'extérieur.}
  \item Ajoutez une méthode \py{set_note} qui n'accepte que des notes entre 0 et 20.
  \item Créez un étudiant nommé \py{"Paul"} avec une note de 15, puis essayez de changer la note
    à 25.
\end{qbloc}
\zone[4.0cm]
\begin{solution}
\begin{lstlisting}
    def set_note(self, nouvelle_note):
        if 0 <= nouvelle_note <= 20:
            self.__note = nouvelle_note
        else:
            print("Note invalide : elle doit etre comprise entre 0 et 20.")

paul = Etudiant("Paul", 15)
paul.set_note(25)         # Note invalide : ...
print(paul.get_note())    # 15 : la note n'a pas change
\end{lstlisting}
\textbf{Variante professionnelle.} Au lieu d'un \py{print}, lever une exception rend l'erreur
impossible à ignorer :\\
\py{raise ValueError("note hors intervalle")}.
\end{solution}

% ---------------------------------------------------------------- B4
\exo{Héritage simple}
\begin{codebloc}
\begin{lstlisting}
class Animal:
    def __init__(self, nom):
        self.__nom = nom

    def get_nom(self):
        return self.__nom

    def faire_son(self):
        return "Son generique"

class Chat(Animal):
    # Méthode faire_son pour retourner "Miaou"
\end{lstlisting}
\end{codebloc}
\begin{qbloc}
  \item Quelle classe est la classe parent ?
    \rep{\py{Animal}. \py{Chat} est la classe fille (ou classe dérivée).}
  \item Que fait l'héritage dans ce cas ?
    \rep{\py{Chat} récupère sans les réécrire le constructeur \py{__init__}, l'attribut privé
    \py{__nom} et la méthode \py{get_nom} d'\py{Animal}. Elle peut en plus redéfinir \py{faire_son}
    pour lui donner son comportement propre.}
  \item Redéfinissez la méthode \py{faire_son} pour retourner \py{"Miaou"}.
  \item Créez un \py{Chat} nommé \py{"Minou"} et affichez son son.
\end{qbloc}
\zone[4.0cm]
\begin{solution}
\begin{lstlisting}
class Chat(Animal):
    def faire_son(self):
        return "Miaou"

minou = Chat("Minou")
print(minou.get_nom())     # Minou  (methode heritee)
print(minou.faire_son())   # Miaou  (methode redefinie)
\end{lstlisting}
\py{Chat} ne définit pas de \py{__init__} : celui d'\py{Animal} est utilisé tel quel. Appeler
\py{faire_son()} sur un \py{Chat} exécute la version de \py{Chat} : c'est le
\textbf{polymorphisme}.
\end{solution}

% ---------------------------------------------------------------- B5
\exo{Héritage avec nouvel attribut}
\begin{codebloc}
\begin{lstlisting}
class Vehicule:
    def __init__(self, marque):
        self.__marque = marque

    def get_marque(self):
        return self.__marque

class Voiture(Vehicule):
    def __init__(self, marque, couleur):
        super().__init__(marque)
        self.__couleur = couleur

    # À compléter : méthode get_couleur
\end{lstlisting}
\end{codebloc}
\begin{qbloc}
  \item Que fait \py{super().__init__(marque)} ?
    \rep{Il appelle le constructeur de la classe mère \py{Vehicule} pour initialiser la partie
    héritée de l'objet, ici l'attribut privé \py{__marque}. Sans cet appel, \py{get_marque}
    échouerait car l'attribut n'aurait jamais été créé.}
  \item Combien d'attributs privés possède la classe \py{Voiture} ?
    \rep{Deux au total dans l'objet : \py{_Vehicule__marque}, hérité et initialisé par \py{super()},
    et \py{_Voiture__couleur}, propre à \py{Voiture}. La classe \py{Voiture} n'en déclare qu'un seul
    elle-même.}
  \item À compléter : ajoutez une méthode \py{get_couleur}.
\end{qbloc}
\zone[4.0cm]
\begin{solution}
\begin{lstlisting}
    def get_couleur(self):
        return self.__couleur

v = Voiture("Renault", "Bleu")
print(v.get_marque())      # Renault
print(v.get_couleur())     # Bleu
\end{lstlisting}
\textcolor{ipred}{\textbf{Piège.}} \py{__marque} est privé \emph{à \py{Vehicule}} : à cause du
\emph{name mangling}, une méthode de \py{Voiture} ne peut pas écrire \py{self.__marque}. Elle doit
passer par \py{self.get_marque()}. C'est exactement le rôle de l'accesseur hérité.
\end{solution}

% ---------------------------------------------------------------- B6
\exo{Encapsulation et mutateur avec validation}
\begin{codebloc}
\begin{lstlisting}
class Livre:
    def __init__(self, titre, prix):
        self.__titre = titre
        self.__prix = prix

    def get_prix(self):
        return self.__prix

    # À compléter : set_prix refuse les prix négatifs
\end{lstlisting}
\end{codebloc}
\begin{qbloc}
  \item Pourquoi le prix est-il un attribut privé ?
    \rep{Parce qu'un prix doit respecter une règle métier : il ne peut pas être négatif. Le rendre
    privé interdit l'écriture directe et oblige à passer par une méthode qui applique cette règle.}
  \item Que fait un mutateur avec validation ?
    \rep{Il vérifie la valeur reçue avant de l'affecter. Si elle est valide, il met l'attribut à
    jour ; sinon il refuse la modification et l'objet conserve son état précédent, qui reste
    cohérent.}
  \item Créez un \py{Livre} intitulé \py{"Python 101"} avec un prix de 30, puis essayez de définir
    un prix de $-10$.
\end{qbloc}
\zone[4.0cm]
\begin{solution}
\begin{lstlisting}
    def set_prix(self, nouveau_prix):
        if nouveau_prix >= 0:
            self.__prix = nouveau_prix
        else:
            print("Prix invalide : un prix ne peut pas etre negatif.")

livre = Livre("Python 101", 30)
livre.set_prix(-10)        # Prix invalide : ...
print(livre.get_prix())    # 30 : le prix est inchange
\end{lstlisting}
\textbf{Remarque.} Le constructeur devrait appliquer la même validation, par exemple en appelant
\py{self.set_prix(prix)} au lieu d'affecter directement : sinon \py{Livre("X", -10)} crée un objet
invalide dès le départ.
\end{solution}

% ============================================================ synthèse (corrigé)
\ifcorrige
\clearpage
\partie{Synthèse}
\begin{center}\small
\begin{tabularx}{\textwidth}{@{}>{\bfseries}p{4cm} X@{}}
\toprule
Notion & Ce qu'il faut retenir \\
\midrule
Classe & Modèle décrivant des données et des comportements : \py{class Voiture:} \\
Objet / instance & Réalisation concrète de la classe : \py{v = Voiture("Toyota", "Bleu")} \\
Constructeur & \py{__init__}, appelé automatiquement, initialise les attributs \\
\py{self} & Référence à l'objet courant ; premier paramètre de toute méthode d'instance \\
Attribut d'instance & Donnée propre à l'objet : \py{self.marque} \\
Méthode & Fonction définie dans la classe, appelée par \py{objet.methode()} \\
Attribut privé & \py{self.__prix} : inaccessible depuis l'extérieur (\emph{name mangling}) \\
Accesseur (\emph{getter}) & \py{get_prix()} : lit un attribut privé \\
Mutateur (\emph{setter}) & \py{set_prix(v)} : écrit un attribut privé après validation \\
Héritage & \py{class Voiture(Vehicule)} : la fille reçoit les membres de la mère \\
\py{super()} & Accède à la classe mère, notamment \py{super().__init__(...)} \\
Redéfinition & La fille réécrit une méthode de la mère (\py{faire_son}) \\
\bottomrule
\end{tabularx}
\end{center}

\begin{tcolorbox}[enhanced,breakable,colback=red!4,colframe=ipred,boxrule=0pt,leftrule=2.5pt,
  arc=0pt,outer arc=0pt,left=5pt,right=5pt,top=3pt,bottom=3pt,fontupper=\small]
\textbf{Erreurs les plus fréquentes.} Oublier \py{self} dans la définition d'une méthode ;
confondre \py{return} (retourne une valeur) et \py{print} (affiche) ; écrire \py{objet.__attribut}
depuis l'extérieur ; oublier \py{super().__init__(...)} dans le constructeur d'une classe fille ;
appeler une méthode sans les parenthèses (\py{v.afficher_details} retourne la méthode, pas son
résultat).
\end{tcolorbox}
\fi

\end{document}
"
\documentclass[10pt,a4paper]{article}

\usepackage[T1]{fontenc}
\usepackage[utf8]{inputenc}
\usepackage{lmodern}
\usepackage[french]{babel}
\usepackage[a4paper,top=0.9cm,bottom=1.1cm,left=1.1cm,right=1.1cm]{geometry}
\usepackage[table]{xcolor}
\usepackage{booktabs,tabularx,array}
\usepackage{enumitem}
\usepackage{needspace}
\usepackage{fancyhdr}
\usepackage{amssymb,amsmath}
\usepackage[most]{tcolorbox}
\usepackage{listings}
\usepackage[strings]{underscore}  % permet d'écrire __init__ sans échappement
\usepackage{hyperref}
\hypersetup{hidelinks}

\providecommand{\modecorrige}{0}
\newif\ifcorrige
\ifnum\modecorrige=1 \corrigetrue \else \corrigefalse \fi

% Palette alignée sur les autres supports du dépôt
\definecolor{ipred}{HTML}{C0392B}
\definecolor{cdcbar}{HTML}{A67C3D}
\definecolor{cdcbg}{HTML}{F4F1EA}
\definecolor{qbar}{HTML}{4A6B5D}
\definecolor{qbg}{HTML}{EDF1EF}
\definecolor{corr}{HTML}{1F5FA8}
\definecolor{corrbg}{HTML}{E8F0FA}
\definecolor{grisfoot}{HTML}{828282}
\definecolor{codebg}{HTML}{F7F7F4}
\definecolor{codekw}{HTML}{8E44AD}
\definecolor{codestr}{HTML}{1F7A4D}
\definecolor{codecom}{HTML}{8A8A8A}

% ---------------------------------------------------------------- code Python
\lstdefinestyle{py}{
  language=Python,
  basicstyle=\ttfamily\scriptsize,
  keywordstyle=\color{codekw}\bfseries,
  stringstyle=\color{codestr},
  commentstyle=\color{codecom}\itshape,
  showstringspaces=false,
  backgroundcolor=\color{codebg},
  frame=single,
  rulecolor=\color{black!25},
  framesep=3pt,
  aboveskip=3pt, belowskip=3pt,
  xleftmargin=2pt, xrightmargin=2pt,
  columns=fullflexible,
  keepspaces=true,
  breaklines=true,
  extendedchars=true,
  literate={é}{{\'e}}1 {è}{{\`e}}1 {ê}{{\^e}}1 {à}{{\`a}}1 {ç}{{\c c}}1
           {ù}{{\`u}}1 {û}{{\^u}}1 {î}{{\^\i}}1 {ô}{{\^o}}1
           {É}{{\'E}}1 {À}{{\`A}}1 {È}{{\`E}}1 {Ç}{{\c C}}1
           {«}{{\guillemotleft}}1 {»}{{\guillemotright}}1,
  morekeywords={self,True,False,None},
}
\lstset{style=py}
\newcommand{\py}[1]{\texttt{#1}}

% --------------------------------------------------------------- encadrés
\newtcolorbox{memo}[1][]{enhanced,breakable,colback=cdcbg,colframe=cdcbar,boxrule=0pt,
  leftrule=2.5pt,arc=0pt,outer arc=0pt,left=5pt,right=5pt,top=3pt,bottom=3pt,
  title={#1},coltitle=cdcbar,fonttitle=\bfseries\footnotesize,attach title to upper={\ :\ },
  fontupper=\footnotesize}

% --------------------------------------------------------------- exercices
\newcounter{exo}
\newcommand{\exo}[1]{%
  \stepcounter{exo}\par\vspace{4pt}
  \ifcorrige\else\noindent\minipage{\textwidth}\fi
  \noindent\textcolor{qbar}{\rule{2.5pt}{10pt}}\hspace{4pt}%
  \textbf{\color{qbar}Exercice \theexo}\ --- \textbf{#1}%
  \par\nobreak\vspace{-3pt}\noindent\textcolor{qbar!35}{\rule{\textwidth}{0.6pt}}%
  \par\nobreak\vspace{1pt}}

% Mise en page de l'énoncé :
%   fiche élève  : code à gauche, questions à droite (gain de place vertical)
%   corrigé      : pleine largeur, l'un sous l'autre (les réponses prennent la place)
\newenvironment{codebloc}
  {\ifcorrige\else\noindent\minipage[t]{0.50\textwidth}\vspace{0pt}\fi}
  {\ifcorrige\else\endminipage\hfill\fi}
\newenvironment{qbloc}
  {\ifcorrige\else\minipage[t]{0.475\textwidth}\vspace{0pt}\fi
   \textbf{\small Questions}\begin{enumerate}[leftmargin=*,nosep,label=\arabic*.]}
  {\end{enumerate}\ifcorrige\else\endminipage\par\fi}

% Variante pleine largeur : exercices sans bloc de code à afficher
\newenvironment{qbloclarge}
  {\par\vspace{2pt}\noindent\textbf{\small Questions}\par\nobreak\begin{enumerate}[leftmargin=*,nosep,label=\arabic*.]}
  {\end{enumerate}}

% Réponse d'une question : rien sur la fiche, texte bleu dans le corrigé
\newcommand{\rep}[1]{\ifcorrige\par\smallskip{\color{corr}\small#1}\par\smallskip\fi}

% Cadre de réponse : présent sur la fiche uniquement.
% \zoneadj retranche la même hauteur à tous les cadres : un seul réglage
% pour ajuster la fiche au nombre de pages voulu.
\newlength{\zoneadj}\setlength{\zoneadj}{0pt}
\newcommand{\zone}[1][4cm]{%
  \ifcorrige\else
    \par\vspace{2pt}\noindent
    \begin{tcolorbox}[colback=white,colframe=black!45,boxrule=0.5pt,arc=1pt,height={\dimexpr#1-\zoneadj\relax},
      left=4pt,right=4pt,top=2pt,bottom=2pt,
      title={\scriptsize\itshape Réponses},coltitle=black!55,fonttitle=\itshape,
      detach title,before upper={\tcbtitle\par}]\end{tcolorbox}%
    \endminipage\par
  \fi}

% Bloc visible seulement dans le corrigé
\newenvironment{solution}
  {\ifcorrige
     \begin{tcolorbox}[enhanced,breakable,colback=corrbg,colframe=corr,boxrule=0.6pt,arc=1pt,
       left=5pt,right=5pt,top=3pt,bottom=3pt,fontupper=\small,
       title={Corrigé},coltitle=corr,fonttitle=\bfseries\small,detach title,
       before upper={\tcbtitle\par\smallskip}]%
   \else
     \setbox0\vbox\bgroup
   \fi}
  {\ifcorrige\end{tcolorbox}\else\egroup\fi}

% Partie
\newcommand{\partie}[1]{%
  \Needspace{7cm}\par\medskip\noindent{\large\bfseries\color{qbar}#1}%
  \par\nobreak\vspace{-2pt}\noindent\textcolor{qbar}{\rule{\textwidth}{1pt}}\par\nobreak\smallskip}

\pagestyle{fancy}
\fancyhf{}
\renewcommand{\headrulewidth}{0pt}
\fancyfoot[L]{\scriptsize\itshape\color{grisfoot}Fiche d'exercices --- POO Python\ifcorrige\ --- CORRIGÉ\fi}
\fancyfoot[R]{\scriptsize\color{grisfoot}\thepage}

\setlist{itemsep=1pt,topsep=2pt,parsep=0pt}
\setlength{\parindent}{0pt}
\setlength{\parskip}{2pt}
\setlength{\columnsep}{8mm}

\begin{document}

\noindent
\begin{minipage}[c]{0.70\textwidth}
{\large\bfseries Fiche d'exercices --- Programmation orientée objet en Python}\\[1pt]
{\small Classes, attributs, méthodes \ $\cdot$\ Encapsulation et héritage \ $\cdot$\ BTS CIEL}
\ifcorrige\\[1pt]{\small\color{corr}\bfseries Version enseignant --- corrigés intégrés}\fi
\end{minipage}\hfill
\begin{minipage}[c]{0.28\textwidth}\small
\ifcorrige\color{corr}corrigé\else Nom :\dotfill\\[3pt] Classe :\dotfill\fi
\end{minipage}

\vspace{2pt}
\begin{memo}[Mémo]
\textbf{Classe} : modèle. \textbf{Objet} (instance) : réalisation concrète de ce modèle.
\textbf{Attribut} : donnée portée par l'objet. \textbf{Méthode} : fonction définie dans la classe.
\py{self} désigne l'objet sur lequel la méthode est appelée ; \py{__init__} est le
\textbf{constructeur}. \py{return} \emph{retourne} une valeur, \py{print} l'\emph{affiche}.
\end{memo}

\partie{Partie A --- Classes, attributs et méthodes}

% ---------------------------------------------------------------- A1
\exo{La classe \py{Voiture}}
\begin{codebloc}
\begin{lstlisting}
class Voiture:
    def __init__(self, marque, couleur):
        self.marque = marque
        self.couleur = couleur

    def afficher_details(self):
        return f"Voiture {self.marque} de couleur {self.couleur}"
\end{lstlisting}
\end{codebloc}
\begin{qbloc}
  \item Que fait la méthode \py{__init__} ?
    \rep{C'est le \textbf{constructeur}. Il est appelé automatiquement à la création de l'objet et
    initialise les attributs d'instance \py{self.marque} et \py{self.couleur} à partir des valeurs
    passées en paramètres.}
  \item Combien d'attributs possède la classe \py{Voiture} ?
    \rep{Deux attributs d'instance : \py{marque} et \py{couleur}.}
  \item Que retourne la méthode \py{afficher_details} ?
    \rep{Une chaîne de caractères (\py{str}) construite par une f-string, par exemple
    \py{"Voiture Toyota de couleur Bleu"}. Elle \emph{retourne} la chaîne, elle ne l'affiche pas.}
  \item Créez une instance de \py{Voiture} de marque \py{"Toyota"} et de couleur \py{"Bleu"},
    appelez \py{afficher_details} et affichez le résultat.
\end{qbloc}
\zone[3.8cm]
\begin{solution}
\begin{lstlisting}
ma_voiture = Voiture("Toyota", "Bleu")
print(ma_voiture.afficher_details())      # Voiture Toyota de couleur Bleu
\end{lstlisting}
\py{print} est nécessaire : sans lui, la chaîne est retournée puis perdue.
\end{solution}

% ---------------------------------------------------------------- A2
\exo{Compléter la classe \py{Chien}}
\begin{codebloc}
\begin{lstlisting}
class Chien:
    def __init__(self, nom, age):
        self.nom = nom
        self.age = age

    # Méthode aboyer qui retourne "Wouf !"
\end{lstlisting}
\end{codebloc}
\begin{qbloc}
  \item Quels sont les attributs de la classe \py{Chien} ?
    \rep{\py{nom} et \py{age}, deux attributs d'instance.}
  \item Comment appelle-t-on la méthode \py{__init__} dans le contexte de la POO ?
    \rep{Le \textbf{constructeur} (ou initialiseur). C'est une \emph{méthode spéciale} ; on ne
    l'appelle jamais directement, elle est déclenchée par \py{Chien("Rex", 3)}.}
  \item Ajoutez une méthode \py{aboyer} qui retourne \py{"Wouf !"}.
  \item Créez un objet \py{Chien} nommé \py{"Rex"}, âgé de 3 ans, et faites-le aboyer.
\end{qbloc}
\zone[3.8cm]
\begin{solution}
\begin{lstlisting}
class Chien:
    def __init__(self, nom, age):
        self.nom = nom
        self.age = age

    def aboyer(self):
        return "Wouf !"

rex = Chien("Rex", 3)
print(rex.aboyer())        # Wouf !
\end{lstlisting}
Le paramètre \py{self} est obligatoire dans la définition, même si \py{aboyer} n'utilise aucun
attribut : il reçoit l'objet sur lequel la méthode est appelée.
\end{solution}

% ---------------------------------------------------------------- A3
\exo{La classe \py{Livre}}
\begin{codebloc}
\begin{lstlisting}
class Livre:
    def __init__(self, titre, auteur):
        self.titre = titre
        self.auteur = auteur

    def description(self):
        return f"{self.titre} par {self.auteur}"
\end{lstlisting}
\end{codebloc}
\begin{qbloc}
  \item Que fait la méthode \py{description} ?
    \rep{Elle retourne une chaîne qui associe le titre et l'auteur du livre, sous la forme
    \py{"titre par auteur"}.}
  \item Combien de paramètres la méthode \py{__init__} attend-elle, en incluant \py{self} ?
    \rep{Trois : \py{self}, \py{titre} et \py{auteur}. À l'appel, seuls deux arguments sont
    fournis : \py{self} est transmis automatiquement par Python.}
  \item Créez un objet \py{Livre} de titre \py{"Le Petit Prince"} et d'auteur
    \py{"Antoine de Saint-Exupéry"}, puis affichez sa description.
\end{qbloc}
\zone[3.8cm]
\begin{solution}
\begin{lstlisting}
livre = Livre("Le Petit Prince", "Antoine de Saint-Exupery")
print(livre.description())      # Le Petit Prince par Antoine de Saint-Exupery
\end{lstlisting}
\end{solution}

% ---------------------------------------------------------------- A4
\exo{La classe \py{Rectangle}}
\begin{codebloc}
\begin{lstlisting}
class Rectangle:
    def __init__(self, longueur, largeur):
        self.longueur = longueur
        self.largeur = largeur

    # Méthode calculer_aire qui retourne l'aire
\end{lstlisting}
\end{codebloc}
\begin{qbloc}
  \item Quels sont les attributs de la classe \py{Rectangle} ?
    \rep{\py{longueur} et \py{largeur}.}
  \item Quelle formule mathématique devez-vous utiliser pour calculer l'aire ?
    \rep{$\text{aire} = \text{longueur} \times \text{largeur}$.}
  \item Ajoutez une méthode \py{calculer_aire} qui retourne l'aire du rectangle.
  \item Créez un rectangle de longueur 5 et de largeur 3, puis affichez son aire.
\end{qbloc}
\zone[8.2cm]
\begin{solution}
\begin{lstlisting}
class Rectangle:
    def __init__(self, longueur, largeur):
        self.longueur = longueur
        self.largeur = largeur

    def calculer_aire(self):
        return self.longueur * self.largeur

r = Rectangle(5, 3)
print(r.calculer_aire())        # 15
\end{lstlisting}
Une méthode peut aussi être ajoutée pour le périmètre :\\
\py{return 2 * (self.longueur + self.largeur)}.
\end{solution}

% ---------------------------------------------------------------- A5
\exo{Création complète d'une classe \py{Etudiant}}
On souhaite modéliser un étudiant par son nom et sa note.
\begin{qbloclarge}
  \item Créez une classe \py{Etudiant} avec un constructeur \py{__init__} prenant en paramètres
    \py{nom} et \py{note}, et deux attributs d'instance \py{self.nom} et \py{self.note}.
  \item Ajoutez une méthode \py{a_reussi} qui retourne \py{True} si la note est supérieure ou égale
    à 10, et \py{False} sinon.
  \item Testez votre classe en créant un étudiant avec une note de 8.
  \item Ajoutez une méthode \py{augmenter_note} qui augmente la note de 2 points.
  \item Créez un étudiant nommé \py{"Alice"} avec une note de 15 et affichez si elle a réussi,
    grâce à la méthode \py{a_reussi}.
\end{qbloclarge}
\zone[8.7cm]
\begin{solution}
\begin{lstlisting}
class Etudiant:
    def __init__(self, nom, note):
        self.nom = nom
        self.note = note

    def a_reussi(self):
        return self.note >= 10

    def augmenter_note(self):
        self.note = self.note + 2

e = Etudiant("Bob", 8)          # 3. test avec une note de 8
print(e.a_reussi())             # False

alice = Etudiant("Alice", 15)   # 5. Alice
print(alice.a_reussi())         # True
\end{lstlisting}
\textbf{Points d'attention.} \py{return self.note >= 10} suffit : la comparaison produit déjà un
booléen, un \py{if ... return True else return False} est inutile. \py{augmenter_note} \emph{modifie}
l'attribut et ne retourne rien : elle agit sur l'état de l'objet.
\end{solution}

\partie{Partie B --- Encapsulation et héritage}

\begin{memo}[Mémo]
\textbf{Encapsulation} : \py{self.__nom} (deux soulignés) déclenche le \emph{name mangling} — Python
renomme l'attribut en \py{_Classe__nom}, donc inaccessible de l'extérieur par son nom d'origine. On y
accède par un \textbf{accesseur} (\py{get_nom}) et on le modifie par un \textbf{mutateur}
(\py{set_nom}), qui peut valider la valeur.
\textbf{Héritage} : \py{class Chat(Animal)} — la classe fille reçoit les attributs et méthodes de la
mère, peut en ajouter ou \textbf{redéfinir} les siennes ; \py{super().__init__(...)} appelle le
constructeur de la mère.
\end{memo}

% ---------------------------------------------------------------- B1
\exo{Encapsulation simple avec attribut privé}
\begin{codebloc}
\begin{lstlisting}
class Personne:
    def __init__(self, nom, age):
        self.__nom = nom
        self.__age = age

    def get_nom(self):
        return self.__nom

    # Méthode get_age()
\end{lstlisting}
\end{codebloc}
\begin{qbloc}
  \item Pourquoi les attributs \py{__nom} et \py{__age} sont-ils privés ?
    \rep{Pour empêcher leur modification directe depuis l'extérieur. L'accès passe par des méthodes,
    ce qui permet de contrôler la lecture et l'écriture et de modifier l'implémentation interne sans
    casser le code qui utilise la classe.}
  \item Comment accède-t-on à l'attribut \py{__nom} ?
    \rep{Par l'accesseur : \py{p.get_nom()}. L'écriture \py{p.__nom} provoque une
    \py{AttributeError}.}
  \item Ajoutez une méthode \py{get_age} pour accéder à l'attribut \py{__age}.
  \item Créez une instance de \py{Personne} de nom \py{"Luc"} et d'âge 25, puis affichez son nom.
\end{qbloc}
\zone[4.0cm]
\begin{solution}
\begin{lstlisting}
    def get_age(self):
        return self.__age

p = Personne("Luc", 25)
print(p.get_nom())        # Luc
print(p.get_age())        # 25
\end{lstlisting}
\end{solution}

% ---------------------------------------------------------------- B2
\exo{Ajouter un mutateur}
\begin{codebloc}
\begin{lstlisting}
class CompteBancaire:
    def __init__(self, titulaire, solde):
        self.__titulaire = titulaire
        self.__solde = solde

    def get_solde(self):
        return self.__solde
\end{lstlisting}
\end{codebloc}
\begin{qbloc}
  \item Que signifie l'encapsulation dans cette classe ?
    \rep{Le titulaire et le solde sont privés : aucun code extérieur ne peut écrire
    \py{compte.__solde = 1000000}. Le solde ne peut être consulté que par \py{get_solde} et modifié
    que par les méthodes que la classe met à disposition.}
  \item Que fait la méthode \py{get_solde} ?
    \rep{Elle retourne la valeur de l'attribut privé \py{__solde} : c'est un accesseur en lecture
    seule.}
  \item Créez un compte bancaire pour \py{"Emma"} avec un solde de 500 et affichez le solde.
  \item Ajoutez une méthode \py{set_solde} pour modifier le solde.
\end{qbloc}
\zone[4.0cm]
\begin{solution}
\begin{lstlisting}
    def set_solde(self, nouveau_solde):
        self.__solde = nouveau_solde

compte = CompteBancaire("Emma", 500)
print(compte.get_solde())       # 500
compte.set_solde(750)
print(compte.get_solde())       # 750
\end{lstlisting}
\textbf{Prolongement.} Un mutateur brut n'apporte aucune sécurité. Une version plus réaliste refuse
un solde négatif, ou expose plutôt des opérations métier : \py{deposer(montant)} et
\py{retirer(montant)} avec vérification de la provision.
\end{solution}

% ---------------------------------------------------------------- B3
\exo{Encapsulation avec validation}
\begin{codebloc}
\begin{lstlisting}
class Etudiant:
    def __init__(self, nom, note):
        self.__nom = nom
        self.__note = note

    def get_note(self):
        return self.__note

    # Méthode set_note : notes entre 0 et 20
\end{lstlisting}
\end{codebloc}
\begin{qbloc}
  \item Pourquoi utiliser un mutateur pour la note ?
    \rep{Pour garantir un \emph{invariant} : une note reste comprise entre 0 et 20. Le mutateur est
    le seul point d'entrée en écriture, il peut donc vérifier la valeur avant de l'enregistrer.}
  \item Que se passe-t-il si on essaie d'accéder directement à \py{__note} ?
    \rep{Python lève une \py{AttributeError} : à cause du \emph{name mangling}, l'attribut s'appelle
    en réalité \py{_Etudiant__note} et le nom \py{__note} n'existe pas depuis l'extérieur.}
  \item Ajoutez une méthode \py{set_note} qui n'accepte que des notes entre 0 et 20.
  \item Créez un étudiant nommé \py{"Paul"} avec une note de 15, puis essayez de changer la note
    à 25.
\end{qbloc}
\zone[4.0cm]
\begin{solution}
\begin{lstlisting}
    def set_note(self, nouvelle_note):
        if 0 <= nouvelle_note <= 20:
            self.__note = nouvelle_note
        else:
            print("Note invalide : elle doit etre comprise entre 0 et 20.")

paul = Etudiant("Paul", 15)
paul.set_note(25)         # Note invalide : ...
print(paul.get_note())    # 15 : la note n'a pas change
\end{lstlisting}
\textbf{Variante professionnelle.} Au lieu d'un \py{print}, lever une exception rend l'erreur
impossible à ignorer :\\
\py{raise ValueError("note hors intervalle")}.
\end{solution}

% ---------------------------------------------------------------- B4
\exo{Héritage simple}
\begin{codebloc}
\begin{lstlisting}
class Animal:
    def __init__(self, nom):
        self.__nom = nom

    def get_nom(self):
        return self.__nom

    def faire_son(self):
        return "Son generique"

class Chat(Animal):
    # Méthode faire_son pour retourner "Miaou"
\end{lstlisting}
\end{codebloc}
\begin{qbloc}
  \item Quelle classe est la classe parent ?
    \rep{\py{Animal}. \py{Chat} est la classe fille (ou classe dérivée).}
  \item Que fait l'héritage dans ce cas ?
    \rep{\py{Chat} récupère sans les réécrire le constructeur \py{__init__}, l'attribut privé
    \py{__nom} et la méthode \py{get_nom} d'\py{Animal}. Elle peut en plus redéfinir \py{faire_son}
    pour lui donner son comportement propre.}
  \item Redéfinissez la méthode \py{faire_son} pour retourner \py{"Miaou"}.
  \item Créez un \py{Chat} nommé \py{"Minou"} et affichez son son.
\end{qbloc}
\zone[4.0cm]
\begin{solution}
\begin{lstlisting}
class Chat(Animal):
    def faire_son(self):
        return "Miaou"

minou = Chat("Minou")
print(minou.get_nom())     # Minou  (methode heritee)
print(minou.faire_son())   # Miaou  (methode redefinie)
\end{lstlisting}
\py{Chat} ne définit pas de \py{__init__} : celui d'\py{Animal} est utilisé tel quel. Appeler
\py{faire_son()} sur un \py{Chat} exécute la version de \py{Chat} : c'est le
\textbf{polymorphisme}.
\end{solution}

% ---------------------------------------------------------------- B5
\exo{Héritage avec nouvel attribut}
\begin{codebloc}
\begin{lstlisting}
class Vehicule:
    def __init__(self, marque):
        self.__marque = marque

    def get_marque(self):
        return self.__marque

class Voiture(Vehicule):
    def __init__(self, marque, couleur):
        super().__init__(marque)
        self.__couleur = couleur

    # À compléter : méthode get_couleur
\end{lstlisting}
\end{codebloc}
\begin{qbloc}
  \item Que fait \py{super().__init__(marque)} ?
    \rep{Il appelle le constructeur de la classe mère \py{Vehicule} pour initialiser la partie
    héritée de l'objet, ici l'attribut privé \py{__marque}. Sans cet appel, \py{get_marque}
    échouerait car l'attribut n'aurait jamais été créé.}
  \item Combien d'attributs privés possède la classe \py{Voiture} ?
    \rep{Deux au total dans l'objet : \py{_Vehicule__marque}, hérité et initialisé par \py{super()},
    et \py{_Voiture__couleur}, propre à \py{Voiture}. La classe \py{Voiture} n'en déclare qu'un seul
    elle-même.}
  \item À compléter : ajoutez une méthode \py{get_couleur}.
\end{qbloc}
\zone[4.0cm]
\begin{solution}
\begin{lstlisting}
    def get_couleur(self):
        return self.__couleur

v = Voiture("Renault", "Bleu")
print(v.get_marque())      # Renault
print(v.get_couleur())     # Bleu
\end{lstlisting}
\textcolor{ipred}{\textbf{Piège.}} \py{__marque} est privé \emph{à \py{Vehicule}} : à cause du
\emph{name mangling}, une méthode de \py{Voiture} ne peut pas écrire \py{self.__marque}. Elle doit
passer par \py{self.get_marque()}. C'est exactement le rôle de l'accesseur hérité.
\end{solution}

% ---------------------------------------------------------------- B6
\exo{Encapsulation et mutateur avec validation}
\begin{codebloc}
\begin{lstlisting}
class Livre:
    def __init__(self, titre, prix):
        self.__titre = titre
        self.__prix = prix

    def get_prix(self):
        return self.__prix

    # À compléter : set_prix refuse les prix négatifs
\end{lstlisting}
\end{codebloc}
\begin{qbloc}
  \item Pourquoi le prix est-il un attribut privé ?
    \rep{Parce qu'un prix doit respecter une règle métier : il ne peut pas être négatif. Le rendre
    privé interdit l'écriture directe et oblige à passer par une méthode qui applique cette règle.}
  \item Que fait un mutateur avec validation ?
    \rep{Il vérifie la valeur reçue avant de l'affecter. Si elle est valide, il met l'attribut à
    jour ; sinon il refuse la modification et l'objet conserve son état précédent, qui reste
    cohérent.}
  \item Créez un \py{Livre} intitulé \py{"Python 101"} avec un prix de 30, puis essayez de définir
    un prix de $-10$.
\end{qbloc}
\zone[4.0cm]
\begin{solution}
\begin{lstlisting}
    def set_prix(self, nouveau_prix):
        if nouveau_prix >= 0:
            self.__prix = nouveau_prix
        else:
            print("Prix invalide : un prix ne peut pas etre negatif.")

livre = Livre("Python 101", 30)
livre.set_prix(-10)        # Prix invalide : ...
print(livre.get_prix())    # 30 : le prix est inchange
\end{lstlisting}
\textbf{Remarque.} Le constructeur devrait appliquer la même validation, par exemple en appelant
\py{self.set_prix(prix)} au lieu d'affecter directement : sinon \py{Livre("X", -10)} crée un objet
invalide dès le départ.
\end{solution}

% ============================================================ synthèse (corrigé)
\ifcorrige
\clearpage
\partie{Synthèse}
\begin{center}\small
\begin{tabularx}{\textwidth}{@{}>{\bfseries}p{4cm} X@{}}
\toprule
Notion & Ce qu'il faut retenir \\
\midrule
Classe & Modèle décrivant des données et des comportements : \py{class Voiture:} \\
Objet / instance & Réalisation concrète de la classe : \py{v = Voiture("Toyota", "Bleu")} \\
Constructeur & \py{__init__}, appelé automatiquement, initialise les attributs \\
\py{self} & Référence à l'objet courant ; premier paramètre de toute méthode d'instance \\
Attribut d'instance & Donnée propre à l'objet : \py{self.marque} \\
Méthode & Fonction définie dans la classe, appelée par \py{objet.methode()} \\
Attribut privé & \py{self.__prix} : inaccessible depuis l'extérieur (\emph{name mangling}) \\
Accesseur (\emph{getter}) & \py{get_prix()} : lit un attribut privé \\
Mutateur (\emph{setter}) & \py{set_prix(v)} : écrit un attribut privé après validation \\
Héritage & \py{class Voiture(Vehicule)} : la fille reçoit les membres de la mère \\
\py{super()} & Accède à la classe mère, notamment \py{super().__init__(...)} \\
Redéfinition & La fille réécrit une méthode de la mère (\py{faire_son}) \\
\bottomrule
\end{tabularx}
\end{center}

\begin{tcolorbox}[enhanced,breakable,colback=red!4,colframe=ipred,boxrule=0pt,leftrule=2.5pt,
  arc=0pt,outer arc=0pt,left=5pt,right=5pt,top=3pt,bottom=3pt,fontupper=\small]
\textbf{Erreurs les plus fréquentes.} Oublier \py{self} dans la définition d'une méthode ;
confondre \py{return} (retourne une valeur) et \py{print} (affiche) ; écrire \py{objet.__attribut}
depuis l'extérieur ; oublier \py{super().__init__(...)} dans le constructeur d'une classe fille ;
appeler une méthode sans les parenthèses (\py{v.afficher_details} retourne la méthode, pas son
résultat).
\end{tcolorbox}
\fi

\end{document}
"
\documentclass[10pt,a4paper]{article}

\usepackage[T1]{fontenc}
\usepackage[utf8]{inputenc}
\usepackage{lmodern}
\usepackage[french]{babel}
\usepackage[a4paper,top=0.9cm,bottom=1.1cm,left=1.1cm,right=1.1cm]{geometry}
\usepackage[table]{xcolor}
\usepackage{booktabs,tabularx,array}
\usepackage{enumitem}
\usepackage{needspace}
\usepackage{fancyhdr}
\usepackage{amssymb,amsmath}
\usepackage[most]{tcolorbox}
\usepackage{listings}
\usepackage[strings]{underscore}  % permet d'écrire __init__ sans échappement
\usepackage{hyperref}
\hypersetup{hidelinks}

\providecommand{\modecorrige}{0}
\newif\ifcorrige
\ifnum\modecorrige=1 \corrigetrue \else \corrigefalse \fi

% Palette alignée sur les autres supports du dépôt
\definecolor{ipred}{HTML}{C0392B}
\definecolor{cdcbar}{HTML}{A67C3D}
\definecolor{cdcbg}{HTML}{F4F1EA}
\definecolor{qbar}{HTML}{4A6B5D}
\definecolor{qbg}{HTML}{EDF1EF}
\definecolor{corr}{HTML}{1F5FA8}
\definecolor{corrbg}{HTML}{E8F0FA}
\definecolor{grisfoot}{HTML}{828282}
\definecolor{codebg}{HTML}{F7F7F4}
\definecolor{codekw}{HTML}{8E44AD}
\definecolor{codestr}{HTML}{1F7A4D}
\definecolor{codecom}{HTML}{8A8A8A}

% ---------------------------------------------------------------- code Python
\lstdefinestyle{py}{
  language=Python,
  basicstyle=\ttfamily\scriptsize,
  keywordstyle=\color{codekw}\bfseries,
  stringstyle=\color{codestr},
  commentstyle=\color{codecom}\itshape,
  showstringspaces=false,
  backgroundcolor=\color{codebg},
  frame=single,
  rulecolor=\color{black!25},
  framesep=3pt,
  aboveskip=3pt, belowskip=3pt,
  xleftmargin=2pt, xrightmargin=2pt,
  columns=fullflexible,
  keepspaces=true,
  breaklines=true,
  extendedchars=true,
  literate={é}{{\'e}}1 {è}{{\`e}}1 {ê}{{\^e}}1 {à}{{\`a}}1 {ç}{{\c c}}1
           {ù}{{\`u}}1 {û}{{\^u}}1 {î}{{\^\i}}1 {ô}{{\^o}}1
           {É}{{\'E}}1 {À}{{\`A}}1 {È}{{\`E}}1 {Ç}{{\c C}}1
           {«}{{\guillemotleft}}1 {»}{{\guillemotright}}1,
  morekeywords={self,True,False,None},
}
\lstset{style=py}
\newcommand{\py}[1]{\texttt{#1}}

% --------------------------------------------------------------- encadrés
\newtcolorbox{memo}[1][]{enhanced,breakable,colback=cdcbg,colframe=cdcbar,boxrule=0pt,
  leftrule=2.5pt,arc=0pt,outer arc=0pt,left=5pt,right=5pt,top=3pt,bottom=3pt,
  title={#1},coltitle=cdcbar,fonttitle=\bfseries\footnotesize,attach title to upper={\ :\ },
  fontupper=\footnotesize}

% --------------------------------------------------------------- exercices
\newcounter{exo}
\newcommand{\exo}[1]{%
  \stepcounter{exo}\par\vspace{4pt}
  \ifcorrige\else\noindent\minipage{\textwidth}\fi
  \noindent\textcolor{qbar}{\rule{2.5pt}{10pt}}\hspace{4pt}%
  \textbf{\color{qbar}Exercice \theexo}\ --- \textbf{#1}%
  \par\nobreak\vspace{-3pt}\noindent\textcolor{qbar!35}{\rule{\textwidth}{0.6pt}}%
  \par\nobreak\vspace{1pt}}

% Mise en page de l'énoncé :
%   fiche élève  : code à gauche, questions à droite (gain de place vertical)
%   corrigé      : pleine largeur, l'un sous l'autre (les réponses prennent la place)
\newenvironment{codebloc}
  {\ifcorrige\else\noindent\minipage[t]{0.50\textwidth}\vspace{0pt}\fi}
  {\ifcorrige\else\endminipage\hfill\fi}
\newenvironment{qbloc}
  {\ifcorrige\else\minipage[t]{0.475\textwidth}\vspace{0pt}\fi
   \textbf{\small Questions}\begin{enumerate}[leftmargin=*,nosep,label=\arabic*.]}
  {\end{enumerate}\ifcorrige\else\endminipage\par\fi}

% Variante pleine largeur : exercices sans bloc de code à afficher
\newenvironment{qbloclarge}
  {\par\vspace{2pt}\noindent\textbf{\small Questions}\par\nobreak\begin{enumerate}[leftmargin=*,nosep,label=\arabic*.]}
  {\end{enumerate}}

% Réponse d'une question : rien sur la fiche, texte bleu dans le corrigé
\newcommand{\rep}[1]{\ifcorrige\par\smallskip{\color{corr}\small#1}\par\smallskip\fi}

% Cadre de réponse : présent sur la fiche uniquement.
% \zoneadj retranche la même hauteur à tous les cadres : un seul réglage
% pour ajuster la fiche au nombre de pages voulu.
\newlength{\zoneadj}\setlength{\zoneadj}{0pt}
\newcommand{\zone}[1][4cm]{%
  \ifcorrige\else
    \par\vspace{2pt}\noindent
    \begin{tcolorbox}[colback=white,colframe=black!45,boxrule=0.5pt,arc=1pt,height={\dimexpr#1-\zoneadj\relax},
      left=4pt,right=4pt,top=2pt,bottom=2pt,
      title={\scriptsize\itshape Réponses},coltitle=black!55,fonttitle=\itshape,
      detach title,before upper={\tcbtitle\par}]\end{tcolorbox}%
    \endminipage\par
  \fi}

% Bloc visible seulement dans le corrigé
\newenvironment{solution}
  {\ifcorrige
     \begin{tcolorbox}[enhanced,breakable,colback=corrbg,colframe=corr,boxrule=0.6pt,arc=1pt,
       left=5pt,right=5pt,top=3pt,bottom=3pt,fontupper=\small,
       title={Corrigé},coltitle=corr,fonttitle=\bfseries\small,detach title,
       before upper={\tcbtitle\par\smallskip}]%
   \else
     \setbox0\vbox\bgroup
   \fi}
  {\ifcorrige\end{tcolorbox}\else\egroup\fi}

% Partie
\newcommand{\partie}[1]{%
  \Needspace{7cm}\par\medskip\noindent{\large\bfseries\color{qbar}#1}%
  \par\nobreak\vspace{-2pt}\noindent\textcolor{qbar}{\rule{\textwidth}{1pt}}\par\nobreak\smallskip}

\pagestyle{fancy}
\fancyhf{}
\renewcommand{\headrulewidth}{0pt}
\fancyfoot[L]{\scriptsize\itshape\color{grisfoot}Fiche d'exercices --- POO Python\ifcorrige\ --- CORRIGÉ\fi}
\fancyfoot[R]{\scriptsize\color{grisfoot}\thepage}

\setlist{itemsep=1pt,topsep=2pt,parsep=0pt}
\setlength{\parindent}{0pt}
\setlength{\parskip}{2pt}
\setlength{\columnsep}{8mm}

\begin{document}

\noindent
\begin{minipage}[c]{0.70\textwidth}
{\large\bfseries Fiche d'exercices --- Programmation orientée objet en Python}\\[1pt]
{\small Classes, attributs, méthodes \ $\cdot$\ Encapsulation et héritage \ $\cdot$\ BTS CIEL}
\ifcorrige\\[1pt]{\small\color{corr}\bfseries Version enseignant --- corrigés intégrés}\fi
\end{minipage}\hfill
\begin{minipage}[c]{0.28\textwidth}\small
\ifcorrige\color{corr}corrigé\else Nom :\dotfill\\[3pt] Classe :\dotfill\fi
\end{minipage}

\vspace{2pt}
\begin{memo}[Mémo]
\textbf{Classe} : modèle. \textbf{Objet} (instance) : réalisation concrète de ce modèle.
\textbf{Attribut} : donnée portée par l'objet. \textbf{Méthode} : fonction définie dans la classe.
\py{self} désigne l'objet sur lequel la méthode est appelée ; \py{__init__} est le
\textbf{constructeur}. \py{return} \emph{retourne} une valeur, \py{print} l'\emph{affiche}.
\end{memo}

\partie{Partie A --- Classes, attributs et méthodes}

% ---------------------------------------------------------------- A1
\exo{La classe \py{Voiture}}
\begin{codebloc}
\begin{lstlisting}
class Voiture:
    def __init__(self, marque, couleur):
        self.marque = marque
        self.couleur = couleur

    def afficher_details(self):
        return f"Voiture {self.marque} de couleur {self.couleur}"
\end{lstlisting}
\end{codebloc}
\begin{qbloc}
  \item Que fait la méthode \py{__init__} ?
    \rep{C'est le \textbf{constructeur}. Il est appelé automatiquement à la création de l'objet et
    initialise les attributs d'instance \py{self.marque} et \py{self.couleur} à partir des valeurs
    passées en paramètres.}
  \item Combien d'attributs possède la classe \py{Voiture} ?
    \rep{Deux attributs d'instance : \py{marque} et \py{couleur}.}
  \item Que retourne la méthode \py{afficher_details} ?
    \rep{Une chaîne de caractères (\py{str}) construite par une f-string, par exemple
    \py{"Voiture Toyota de couleur Bleu"}. Elle \emph{retourne} la chaîne, elle ne l'affiche pas.}
  \item Créez une instance de \py{Voiture} de marque \py{"Toyota"} et de couleur \py{"Bleu"},
    appelez \py{afficher_details} et affichez le résultat.
\end{qbloc}
\zone[3.8cm]
\begin{solution}
\begin{lstlisting}
ma_voiture = Voiture("Toyota", "Bleu")
print(ma_voiture.afficher_details())      # Voiture Toyota de couleur Bleu
\end{lstlisting}
\py{print} est nécessaire : sans lui, la chaîne est retournée puis perdue.
\end{solution}

% ---------------------------------------------------------------- A2
\exo{Compléter la classe \py{Chien}}
\begin{codebloc}
\begin{lstlisting}
class Chien:
    def __init__(self, nom, age):
        self.nom = nom
        self.age = age

    # Méthode aboyer qui retourne "Wouf !"
\end{lstlisting}
\end{codebloc}
\begin{qbloc}
  \item Quels sont les attributs de la classe \py{Chien} ?
    \rep{\py{nom} et \py{age}, deux attributs d'instance.}
  \item Comment appelle-t-on la méthode \py{__init__} dans le contexte de la POO ?
    \rep{Le \textbf{constructeur} (ou initialiseur). C'est une \emph{méthode spéciale} ; on ne
    l'appelle jamais directement, elle est déclenchée par \py{Chien("Rex", 3)}.}
  \item Ajoutez une méthode \py{aboyer} qui retourne \py{"Wouf !"}.
  \item Créez un objet \py{Chien} nommé \py{"Rex"}, âgé de 3 ans, et faites-le aboyer.
\end{qbloc}
\zone[3.8cm]
\begin{solution}
\begin{lstlisting}
class Chien:
    def __init__(self, nom, age):
        self.nom = nom
        self.age = age

    def aboyer(self):
        return "Wouf !"

rex = Chien("Rex", 3)
print(rex.aboyer())        # Wouf !
\end{lstlisting}
Le paramètre \py{self} est obligatoire dans la définition, même si \py{aboyer} n'utilise aucun
attribut : il reçoit l'objet sur lequel la méthode est appelée.
\end{solution}

% ---------------------------------------------------------------- A3
\exo{La classe \py{Livre}}
\begin{codebloc}
\begin{lstlisting}
class Livre:
    def __init__(self, titre, auteur):
        self.titre = titre
        self.auteur = auteur

    def description(self):
        return f"{self.titre} par {self.auteur}"
\end{lstlisting}
\end{codebloc}
\begin{qbloc}
  \item Que fait la méthode \py{description} ?
    \rep{Elle retourne une chaîne qui associe le titre et l'auteur du livre, sous la forme
    \py{"titre par auteur"}.}
  \item Combien de paramètres la méthode \py{__init__} attend-elle, en incluant \py{self} ?
    \rep{Trois : \py{self}, \py{titre} et \py{auteur}. À l'appel, seuls deux arguments sont
    fournis : \py{self} est transmis automatiquement par Python.}
  \item Créez un objet \py{Livre} de titre \py{"Le Petit Prince"} et d'auteur
    \py{"Antoine de Saint-Exupéry"}, puis affichez sa description.
\end{qbloc}
\zone[3.8cm]
\begin{solution}
\begin{lstlisting}
livre = Livre("Le Petit Prince", "Antoine de Saint-Exupery")
print(livre.description())      # Le Petit Prince par Antoine de Saint-Exupery
\end{lstlisting}
\end{solution}

% ---------------------------------------------------------------- A4
\exo{La classe \py{Rectangle}}
\begin{codebloc}
\begin{lstlisting}
class Rectangle:
    def __init__(self, longueur, largeur):
        self.longueur = longueur
        self.largeur = largeur

    # Méthode calculer_aire qui retourne l'aire
\end{lstlisting}
\end{codebloc}
\begin{qbloc}
  \item Quels sont les attributs de la classe \py{Rectangle} ?
    \rep{\py{longueur} et \py{largeur}.}
  \item Quelle formule mathématique devez-vous utiliser pour calculer l'aire ?
    \rep{$\text{aire} = \text{longueur} \times \text{largeur}$.}
  \item Ajoutez une méthode \py{calculer_aire} qui retourne l'aire du rectangle.
  \item Créez un rectangle de longueur 5 et de largeur 3, puis affichez son aire.
\end{qbloc}
\zone[8.2cm]
\begin{solution}
\begin{lstlisting}
class Rectangle:
    def __init__(self, longueur, largeur):
        self.longueur = longueur
        self.largeur = largeur

    def calculer_aire(self):
        return self.longueur * self.largeur

r = Rectangle(5, 3)
print(r.calculer_aire())        # 15
\end{lstlisting}
Une méthode peut aussi être ajoutée pour le périmètre :\\
\py{return 2 * (self.longueur + self.largeur)}.
\end{solution}

% ---------------------------------------------------------------- A5
\exo{Création complète d'une classe \py{Etudiant}}
On souhaite modéliser un étudiant par son nom et sa note.
\begin{qbloclarge}
  \item Créez une classe \py{Etudiant} avec un constructeur \py{__init__} prenant en paramètres
    \py{nom} et \py{note}, et deux attributs d'instance \py{self.nom} et \py{self.note}.
  \item Ajoutez une méthode \py{a_reussi} qui retourne \py{True} si la note est supérieure ou égale
    à 10, et \py{False} sinon.
  \item Testez votre classe en créant un étudiant avec une note de 8.
  \item Ajoutez une méthode \py{augmenter_note} qui augmente la note de 2 points.
  \item Créez un étudiant nommé \py{"Alice"} avec une note de 15 et affichez si elle a réussi,
    grâce à la méthode \py{a_reussi}.
\end{qbloclarge}
\zone[8.7cm]
\begin{solution}
\begin{lstlisting}
class Etudiant:
    def __init__(self, nom, note):
        self.nom = nom
        self.note = note

    def a_reussi(self):
        return self.note >= 10

    def augmenter_note(self):
        self.note = self.note + 2

e = Etudiant("Bob", 8)          # 3. test avec une note de 8
print(e.a_reussi())             # False

alice = Etudiant("Alice", 15)   # 5. Alice
print(alice.a_reussi())         # True
\end{lstlisting}
\textbf{Points d'attention.} \py{return self.note >= 10} suffit : la comparaison produit déjà un
booléen, un \py{if ... return True else return False} est inutile. \py{augmenter_note} \emph{modifie}
l'attribut et ne retourne rien : elle agit sur l'état de l'objet.
\end{solution}

\partie{Partie B --- Encapsulation et héritage}

\begin{memo}[Mémo]
\textbf{Encapsulation} : \py{self.__nom} (deux soulignés) déclenche le \emph{name mangling} — Python
renomme l'attribut en \py{_Classe__nom}, donc inaccessible de l'extérieur par son nom d'origine. On y
accède par un \textbf{accesseur} (\py{get_nom}) et on le modifie par un \textbf{mutateur}
(\py{set_nom}), qui peut valider la valeur.
\textbf{Héritage} : \py{class Chat(Animal)} — la classe fille reçoit les attributs et méthodes de la
mère, peut en ajouter ou \textbf{redéfinir} les siennes ; \py{super().__init__(...)} appelle le
constructeur de la mère.
\end{memo}

% ---------------------------------------------------------------- B1
\exo{Encapsulation simple avec attribut privé}
\begin{codebloc}
\begin{lstlisting}
class Personne:
    def __init__(self, nom, age):
        self.__nom = nom
        self.__age = age

    def get_nom(self):
        return self.__nom

    # Méthode get_age()
\end{lstlisting}
\end{codebloc}
\begin{qbloc}
  \item Pourquoi les attributs \py{__nom} et \py{__age} sont-ils privés ?
    \rep{Pour empêcher leur modification directe depuis l'extérieur. L'accès passe par des méthodes,
    ce qui permet de contrôler la lecture et l'écriture et de modifier l'implémentation interne sans
    casser le code qui utilise la classe.}
  \item Comment accède-t-on à l'attribut \py{__nom} ?
    \rep{Par l'accesseur : \py{p.get_nom()}. L'écriture \py{p.__nom} provoque une
    \py{AttributeError}.}
  \item Ajoutez une méthode \py{get_age} pour accéder à l'attribut \py{__age}.
  \item Créez une instance de \py{Personne} de nom \py{"Luc"} et d'âge 25, puis affichez son nom.
\end{qbloc}
\zone[4.0cm]
\begin{solution}
\begin{lstlisting}
    def get_age(self):
        return self.__age

p = Personne("Luc", 25)
print(p.get_nom())        # Luc
print(p.get_age())        # 25
\end{lstlisting}
\end{solution}

% ---------------------------------------------------------------- B2
\exo{Ajouter un mutateur}
\begin{codebloc}
\begin{lstlisting}
class CompteBancaire:
    def __init__(self, titulaire, solde):
        self.__titulaire = titulaire
        self.__solde = solde

    def get_solde(self):
        return self.__solde
\end{lstlisting}
\end{codebloc}
\begin{qbloc}
  \item Que signifie l'encapsulation dans cette classe ?
    \rep{Le titulaire et le solde sont privés : aucun code extérieur ne peut écrire
    \py{compte.__solde = 1000000}. Le solde ne peut être consulté que par \py{get_solde} et modifié
    que par les méthodes que la classe met à disposition.}
  \item Que fait la méthode \py{get_solde} ?
    \rep{Elle retourne la valeur de l'attribut privé \py{__solde} : c'est un accesseur en lecture
    seule.}
  \item Créez un compte bancaire pour \py{"Emma"} avec un solde de 500 et affichez le solde.
  \item Ajoutez une méthode \py{set_solde} pour modifier le solde.
\end{qbloc}
\zone[4.0cm]
\begin{solution}
\begin{lstlisting}
    def set_solde(self, nouveau_solde):
        self.__solde = nouveau_solde

compte = CompteBancaire("Emma", 500)
print(compte.get_solde())       # 500
compte.set_solde(750)
print(compte.get_solde())       # 750
\end{lstlisting}
\textbf{Prolongement.} Un mutateur brut n'apporte aucune sécurité. Une version plus réaliste refuse
un solde négatif, ou expose plutôt des opérations métier : \py{deposer(montant)} et
\py{retirer(montant)} avec vérification de la provision.
\end{solution}

% ---------------------------------------------------------------- B3
\exo{Encapsulation avec validation}
\begin{codebloc}
\begin{lstlisting}
class Etudiant:
    def __init__(self, nom, note):
        self.__nom = nom
        self.__note = note

    def get_note(self):
        return self.__note

    # Méthode set_note : notes entre 0 et 20
\end{lstlisting}
\end{codebloc}
\begin{qbloc}
  \item Pourquoi utiliser un mutateur pour la note ?
    \rep{Pour garantir un \emph{invariant} : une note reste comprise entre 0 et 20. Le mutateur est
    le seul point d'entrée en écriture, il peut donc vérifier la valeur avant de l'enregistrer.}
  \item Que se passe-t-il si on essaie d'accéder directement à \py{__note} ?
    \rep{Python lève une \py{AttributeError} : à cause du \emph{name mangling}, l'attribut s'appelle
    en réalité \py{_Etudiant__note} et le nom \py{__note} n'existe pas depuis l'extérieur.}
  \item Ajoutez une méthode \py{set_note} qui n'accepte que des notes entre 0 et 20.
  \item Créez un étudiant nommé \py{"Paul"} avec une note de 15, puis essayez de changer la note
    à 25.
\end{qbloc}
\zone[4.0cm]
\begin{solution}
\begin{lstlisting}
    def set_note(self, nouvelle_note):
        if 0 <= nouvelle_note <= 20:
            self.__note = nouvelle_note
        else:
            print("Note invalide : elle doit etre comprise entre 0 et 20.")

paul = Etudiant("Paul", 15)
paul.set_note(25)         # Note invalide : ...
print(paul.get_note())    # 15 : la note n'a pas change
\end{lstlisting}
\textbf{Variante professionnelle.} Au lieu d'un \py{print}, lever une exception rend l'erreur
impossible à ignorer :\\
\py{raise ValueError("note hors intervalle")}.
\end{solution}

% ---------------------------------------------------------------- B4
\exo{Héritage simple}
\begin{codebloc}
\begin{lstlisting}
class Animal:
    def __init__(self, nom):
        self.__nom = nom

    def get_nom(self):
        return self.__nom

    def faire_son(self):
        return "Son generique"

class Chat(Animal):
    # Méthode faire_son pour retourner "Miaou"
\end{lstlisting}
\end{codebloc}
\begin{qbloc}
  \item Quelle classe est la classe parent ?
    \rep{\py{Animal}. \py{Chat} est la classe fille (ou classe dérivée).}
  \item Que fait l'héritage dans ce cas ?
    \rep{\py{Chat} récupère sans les réécrire le constructeur \py{__init__}, l'attribut privé
    \py{__nom} et la méthode \py{get_nom} d'\py{Animal}. Elle peut en plus redéfinir \py{faire_son}
    pour lui donner son comportement propre.}
  \item Redéfinissez la méthode \py{faire_son} pour retourner \py{"Miaou"}.
  \item Créez un \py{Chat} nommé \py{"Minou"} et affichez son son.
\end{qbloc}
\zone[4.0cm]
\begin{solution}
\begin{lstlisting}
class Chat(Animal):
    def faire_son(self):
        return "Miaou"

minou = Chat("Minou")
print(minou.get_nom())     # Minou  (methode heritee)
print(minou.faire_son())   # Miaou  (methode redefinie)
\end{lstlisting}
\py{Chat} ne définit pas de \py{__init__} : celui d'\py{Animal} est utilisé tel quel. Appeler
\py{faire_son()} sur un \py{Chat} exécute la version de \py{Chat} : c'est le
\textbf{polymorphisme}.
\end{solution}

% ---------------------------------------------------------------- B5
\exo{Héritage avec nouvel attribut}
\begin{codebloc}
\begin{lstlisting}
class Vehicule:
    def __init__(self, marque):
        self.__marque = marque

    def get_marque(self):
        return self.__marque

class Voiture(Vehicule):
    def __init__(self, marque, couleur):
        super().__init__(marque)
        self.__couleur = couleur

    # À compléter : méthode get_couleur
\end{lstlisting}
\end{codebloc}
\begin{qbloc}
  \item Que fait \py{super().__init__(marque)} ?
    \rep{Il appelle le constructeur de la classe mère \py{Vehicule} pour initialiser la partie
    héritée de l'objet, ici l'attribut privé \py{__marque}. Sans cet appel, \py{get_marque}
    échouerait car l'attribut n'aurait jamais été créé.}
  \item Combien d'attributs privés possède la classe \py{Voiture} ?
    \rep{Deux au total dans l'objet : \py{_Vehicule__marque}, hérité et initialisé par \py{super()},
    et \py{_Voiture__couleur}, propre à \py{Voiture}. La classe \py{Voiture} n'en déclare qu'un seul
    elle-même.}
  \item À compléter : ajoutez une méthode \py{get_couleur}.
\end{qbloc}
\zone[4.0cm]
\begin{solution}
\begin{lstlisting}
    def get_couleur(self):
        return self.__couleur

v = Voiture("Renault", "Bleu")
print(v.get_marque())      # Renault
print(v.get_couleur())     # Bleu
\end{lstlisting}
\textcolor{ipred}{\textbf{Piège.}} \py{__marque} est privé \emph{à \py{Vehicule}} : à cause du
\emph{name mangling}, une méthode de \py{Voiture} ne peut pas écrire \py{self.__marque}. Elle doit
passer par \py{self.get_marque()}. C'est exactement le rôle de l'accesseur hérité.
\end{solution}

% ---------------------------------------------------------------- B6
\exo{Encapsulation et mutateur avec validation}
\begin{codebloc}
\begin{lstlisting}
class Livre:
    def __init__(self, titre, prix):
        self.__titre = titre
        self.__prix = prix

    def get_prix(self):
        return self.__prix

    # À compléter : set_prix refuse les prix négatifs
\end{lstlisting}
\end{codebloc}
\begin{qbloc}
  \item Pourquoi le prix est-il un attribut privé ?
    \rep{Parce qu'un prix doit respecter une règle métier : il ne peut pas être négatif. Le rendre
    privé interdit l'écriture directe et oblige à passer par une méthode qui applique cette règle.}
  \item Que fait un mutateur avec validation ?
    \rep{Il vérifie la valeur reçue avant de l'affecter. Si elle est valide, il met l'attribut à
    jour ; sinon il refuse la modification et l'objet conserve son état précédent, qui reste
    cohérent.}
  \item Créez un \py{Livre} intitulé \py{"Python 101"} avec un prix de 30, puis essayez de définir
    un prix de $-10$.
\end{qbloc}
\zone[4.0cm]
\begin{solution}
\begin{lstlisting}
    def set_prix(self, nouveau_prix):
        if nouveau_prix >= 0:
            self.__prix = nouveau_prix
        else:
            print("Prix invalide : un prix ne peut pas etre negatif.")

livre = Livre("Python 101", 30)
livre.set_prix(-10)        # Prix invalide : ...
print(livre.get_prix())    # 30 : le prix est inchange
\end{lstlisting}
\textbf{Remarque.} Le constructeur devrait appliquer la même validation, par exemple en appelant
\py{self.set_prix(prix)} au lieu d'affecter directement : sinon \py{Livre("X", -10)} crée un objet
invalide dès le départ.
\end{solution}

% ============================================================ synthèse (corrigé)
\ifcorrige
\clearpage
\partie{Synthèse}
\begin{center}\small
\begin{tabularx}{\textwidth}{@{}>{\bfseries}p{4cm} X@{}}
\toprule
Notion & Ce qu'il faut retenir \\
\midrule
Classe & Modèle décrivant des données et des comportements : \py{class Voiture:} \\
Objet / instance & Réalisation concrète de la classe : \py{v = Voiture("Toyota", "Bleu")} \\
Constructeur & \py{__init__}, appelé automatiquement, initialise les attributs \\
\py{self} & Référence à l'objet courant ; premier paramètre de toute méthode d'instance \\
Attribut d'instance & Donnée propre à l'objet : \py{self.marque} \\
Méthode & Fonction définie dans la classe, appelée par \py{objet.methode()} \\
Attribut privé & \py{self.__prix} : inaccessible depuis l'extérieur (\emph{name mangling}) \\
Accesseur (\emph{getter}) & \py{get_prix()} : lit un attribut privé \\
Mutateur (\emph{setter}) & \py{set_prix(v)} : écrit un attribut privé après validation \\
Héritage & \py{class Voiture(Vehicule)} : la fille reçoit les membres de la mère \\
\py{super()} & Accède à la classe mère, notamment \py{super().__init__(...)} \\
Redéfinition & La fille réécrit une méthode de la mère (\py{faire_son}) \\
\bottomrule
\end{tabularx}
\end{center}

\begin{tcolorbox}[enhanced,breakable,colback=red!4,colframe=ipred,boxrule=0pt,leftrule=2.5pt,
  arc=0pt,outer arc=0pt,left=5pt,right=5pt,top=3pt,bottom=3pt,fontupper=\small]
\textbf{Erreurs les plus fréquentes.} Oublier \py{self} dans la définition d'une méthode ;
confondre \py{return} (retourne une valeur) et \py{print} (affiche) ; écrire \py{objet.__attribut}
depuis l'extérieur ; oublier \py{super().__init__(...)} dans le constructeur d'une classe fille ;
appeler une méthode sans les parenthèses (\py{v.afficher_details} retourne la méthode, pas son
résultat).
\end{tcolorbox}
\fi

\end{document}
